% !TEX root = ../thesis_book_0421052099.tex
\chapter{Evaluation: Setup, Results, and Error Analysis}\label{ch:evaluation}

The evaluation is reported in three parts: the experimental design fixed before
any test question was answered (\Cref{sec:setup}), the measurements it produced
(\Cref{sec:results}), and a hand reading of every failure that remained
(\Cref{sec:erroranalysis}).

\section{Experimental Setup}\label{sec:setup}

This part of the evaluation fixes everything that had to be decided before any
test question was answered: which questions, which systems, which numbers, and
which rules for reading them. It records the decisions that went the other way
as well as the ones that stood, each where it belongs rather than collected into
a confessional at the end.

% \allowbreak offers the path separators as break points; without them each path
% is one unbreakable block and the line before it stretches to fit.
The aggregate figures quoted in this chapter are regenerated by
\textsf{scripts/\allowbreak build\_thesis\_numbers.py} into
\textsf{results/\allowbreak phase4/\allowbreak thesis\_numbers.json} from the
scoring, groundedness, judge, and census artifacts. Anything read instead from a
run record or a census sheet is cited to it there. Two \emph{reported} figures
cannot be recomputed from this repository and say so where they appear: the API
spend of \Cref{sec:test-sets}, a billing record, and the second through fourth
determinism rounds of \Cref{sec:backbone}, read from console tables that were
never committed. Most tables are hand copies checked against that JSON rather
than emitted from it. That is the weaker guarantee.

\subsection{Mapping Experiments to Research Questions}\label{sec:experiment-map}

Each research question of \Cref{sec:rqs} is answered by a distinct experiment
rather than by one table read three ways. \textbf{RQ1} is answered by the main
matrix: five systems on two datasets, scored with Hits@1 and F1, broken down by
hop stratum, with paired significance tests
(\Cref{sec:webqsp-results,sec:cwq-results,sec:hop-results}). The claim about
\emph{multi-hop} accuracy rests on the per-stratum breakdown, not on the
aggregate. \textbf{RQ2} is answered by three measurements that disagree with one
another. That disagreement is itself the finding. The three are the two-tier
groundedness metric (\Cref{sec:groundedness-results}), the precision and F1
columns of the main matrix, and the verifier ablation (\Cref{sec:ablations}).
\textbf{RQ3} is answered by four single-deviation ablations on stratified
half-samples, each compared against the unmodified system restricted to the same
questions (\Cref{sec:ablation-conditions,sec:ablations}).

\subsection{Pre-Specification and the No-Tuning
Policy}\label{sec:pre-registration}

One rule governed everything from the moment the test files were built:
\emph{nothing gets tuned}. No prompt was edited, no configuration value changed,
and no threshold adjusted after test data was touched. Defects discovered
mid-matrix were documented rather than fixed. Four decisions were fixed in
writing before the measurements they govern: the frozen configuration $\alpha =
0.7$, $\tau = 0.20$ with claim verification enabled, tuned on the development
set (\Cref{sec:hyperparams}) and never revisited; the sample size of $400$
questions per dataset, named as the fallback before the API budget was known
(\Cref{sec:test-sets}); the baseline certification threshold, that a retrieval
baseline hedges with the gold answer visible in its own retrieved facts on fewer
than 15\% of its hedges (\Cref{sec:baselines}); and the judge validation
threshold of Cohen's $\kappa \geq 0.7$ against blind human annotation
(\Cref{sec:judge}). The last was missed, by $0.0005$, and is reported as missed.
A threshold that is only pre-specified when it is met is not pre-specified at
all. Nothing was filed with a public registry, so the claim throughout is
pre-specification rather than pre-registration.

The single code change Phase~4 required --- a \textsf{use\_planner} flag to
disable decomposition --- was certified as a no-op on the frozen path by
re-running an already-answered development question and confirming that every
language-model call was a cache hit. The reference condition in the ablation
table is therefore provably the same system that produced the main matrix.

\label{sec:metrics-not-run}%
Path fidelity --- agreement between the relation chain the system traversed and
the chain the benchmark's gold SPARQL query specifies --- was the third
criterion named in the approved proposal. It is not reported. It requires the
gold relation chains. The RoG distribution (\Cref{sec:source-data}) publishes
name-keyed triples without the originating SPARQL. So the chains would have to
be reconstructed from the upstream releases and re-aligned to this environment's
relation vocabulary. That was outside the project's remaining budget.
\Cref{sec:future-path-fidelity} treats it as future work.

\section{Evaluation Datasets and Test Sets}\label{sec:test-sets}

\subsection{WebQSP and ComplexWebQuestions}

Two benchmarks are used, in the RoG distribution~\cite{luo2024reasoning}
whose\index{WebQSP}\index{ComplexWebQuestions}\index{dataset} per-question
subgraphs also form the knowledge environment (\Cref{sec:source-data}).
\textbf{WebQSP}~\cite{yih2016value} is the semantic parses of WebQuestions:
natural questions from search logs, mostly one or two hops, with multi-valued
answers common. \textbf{ComplexWebQuestions} (CWQ)~\cite{talmor2018web} is built
by composing WebQSP questions with additional constraints through SPARQL
templates, giving deeper chains, conjunctions, superlatives, and comparisons ---
and, as \Cref{sec:gold-quality} documents, a higher rate of machine-generated
label defects.

The distinction that matters operationally is visible in the gold answer sets.
Counted as the metrics count them --- over distinct answer strings, since
\Cref{sec:kgqa} defines precision and recall over sets --- WebQSP averages
$7.15$ gold entities per question. But its median is $1.5$. Exactly $200$ of its
$400$ questions carry a single gold answer. The mean is pulled up by a long tail
whose largest question lists $237$. CWQ averages $1.75$ with a median of
$1$.\footnote{The distinction matters. The raw answer lists repeat a string on
$4$ WebQSP and $6$ CWQ questions, all of them in the tail. Counting the repeats
takes the WebQSP mean to $7.98$ and its largest question to $382$. But the
scorer deduplicates before matching, so nothing scored in this thesis ever sees
them. Quoting the raw counts would describe a population that no reported number
uses.} So the two benchmarks differ less in the typical question than the means
suggest, and more in the tail. A metric that rewards one correct entity behaves
very differently on a question with hundreds of gold answers than on one with a
single answer. That tail is one reason this thesis reports F1 alongside Hits@1
(\Cref{sec:metrics}).

\subsection{Stratified Sampling, Seeds, and Published Question IDs}

The evaluation uses only the \emph{test} splits. The development set of
\Cref{sec:hyperparams} came from the train and validation splits, and the
test-set builder asserts that the two pools do not overlap.

Sample size was a budget decision. The full test splits hold $1{,}628$ and
$3{,}531$ questions. At the measured token rates, running five systems over both
projected to roughly \$45--55, against a stated ceiling of \$15--20. So the
evaluation fell back to the pre-specified alternative: a stratified sample of
$400$ questions per dataset (\Cref{tab:testsets}). At that size the bootstrap
confidence intervals on Hits@1 are roughly $\pm 5$ points. The arithmetic was
written down before any run fired.

In the end the whole benchmark, ablations included, cost \$11.13 in API spend.
That figure comes from the provider's billing record, not from the run
artifacts. The reproducible proxy is the token and call accounting of
\Cref{sec:efficiency}, which every run record carries per question.

Sampling is proportional within hop strata, at seed~$42$. The resulting question
sets are committed as immutable artifacts
(\textsf{results/phase4/test\_webqsp.json},
\textsf{results/phase4/test\_cwq.json}), and \Cref{app:samples} lists every
question identifier. Once the first system ran, the files never changed again.

\paragraph{A near miss worth recording.} The first build read the training
shards rather than the test shards. The development-overlap guard caught it ---
the contamination guard, built for hygiene, is what found the contamination. No
system had run against those files. They were deleted and rebuilt, and the
rebuilt pools report $1{,}628$ and $3{,}531$ with \textsf{skipped=\{\}}.
\Cref{app:testset-nearmiss} records the diagnosis. The lesson is cheap: an
overlap assertion belongs in every benchmark builder.

\subsection{Hop-Count Stratification and the Accuracy Ceiling}

The stratification variable\index{hop stratum} is the length of the shortest
path, in \emph{this} environment, from any topic entity of the question to any
gold answer entity, capped at four raw hops. The same machinery certified the
development set. So mediator nodes count the same way throughout, as the two
edges they are (\Cref{sec:tool-adjacency}). Each question falls into one of four
strata: \textsf{h1}, \textsf{h2}, \textsf{h3plus}, or \textsf{unreachable}.

\begin{table}[!tb]
  \begin{center}
    \caption[The two certified test samples]{The two certified test samples. The
    unreachable stratum is kept on purpose. The ceiling is the share of the
    $400$ questions with at least one gold answer reachable within four hops. No
    result in \Cref{sec:results} can score above it on Hits@1.}
    \label{tab:testsets}
    \begin{tabular}{|l|r|r|r|r|r|r|}
      \hline
      \textbf{Dataset} & \textbf{n} & \textbf{h1} & \textbf{h2} &
      \textbf{h3plus} & \textbf{unreach.} & \textbf{Ceiling} \\
      \hline
      WebQSP & 400 & 256 & 128 & 4  & 12 & 97.0\% \\
      CWQ    & 400 & 137 & 211 & 49 & 3  & 99.2\% \\
      \hline
    \end{tabular}
  \end{center}
\end{table}

Two facts in \Cref{tab:testsets} come back again and again in the results.
WebQSP is $64\%$ one-hop, and its \textsf{h3plus} tail holds four questions.
Four is too few to interpret, and the tail is labelled as such wherever it
appears. CWQ is mostly two-hop, with a real deep tail of $49$ questions. Because
the samples keep the stratum proportions of their parent pools, per-stratum
numbers describe the benchmark and not just the draw.

The ceilings in the table differ from the ones the validation gate reported. The
difference is explained here rather than left for a reader to reconcile.
\Cref{tab:coverage} measured the full test splits and found $97.3\%$ and
$99.7\%$ of questions reachable. The table above measures the same quantity over
the $400$ questions actually evaluated and finds $97.0\%$ and $99.2\%$. Neither
figure corrects the other. The first says what the environment could support
across the whole benchmark. The second says what any system in
\Cref{sec:results} could have scored, and every accuracy in this thesis is read
against it. The gap between the two is sampling. It stays under a point on both
datasets because the draw kept the stratum proportions
(\Cref{app:ceiling-reconciliation}).

The \textsf{unreachable} stratum stays in. Dropping it would raise every
system's score by the same amount and misdescribe the environment. Keeping it
preserves an honest ceiling. It also turns these questions into the sharpest
available probe of parametric leakage, as \Cref{sec:groundedness-results} shows.
Any ``hit'' on a question with no graph path to gold is an assertion the
environment cannot support.

\subsection{Gold-Answer Quality Control}\label{sec:gold-quality}

Both benchmarks carry label defects, CWQ more. Its questions\index{gold noise}
and answers are generated by SPARQL templates over the knowledge base rather
than written by hand. A question whose gold answer is wrong is not a system
failure. Counting it as one corrupts both the accuracy tables and the error
taxonomy. Some detection procedure was therefore necessary --- though a naive
one is worse than none, as the results below show.

The detection principle is \textbf{independent multi-system consensus against
gold}: when at least three of the five systems of the main matrix assert the
same non-gold answer, the label is worth examining. The voters are those five
systems, spanning two knowledge sources and four retrieval paradigms, rather
than the four ablation conditions, whose agreement would describe an
architecture rather than a label. \Cref{app:goldpass} gives the procedure, the
three verdicts, the automatic triage that cleared $15$ of $89$ rows on WebQSP
and $2$ of $59$ on CWQ, and the two \emph{opposite} evidence signatures that
make the adjudication reproducible --- a world-fact error requiring mixed
consensus, an annotation-pipeline error showing graph-only consensus. One row is
emitted per (question, consensus answer) pair, so rows and questions are
separate units throughout and question counts are always deduplicated.

\paragraph{Census-Based Exclusions and the Dual-Reporting Policy}

\begin{table}[!tb]
  \begin{center}
    \caption[Gold-label adjudication]{Gold-label adjudication. Rows and distinct
questions are separate units because the pass emits one row per consensus
answer. Exclusions are the union of \textsf{gold\_wrong} and
\textsf{ambiguous\_question} questions.}
    \label{tab:goldnoise}
    \begin{tabular}{|l|r|r|r|r|r|}
      \hline
      \textbf{Dataset} & \textbf{Rows} & \textbf{Questions} &
      \textbf{gold\_wrong} & \textbf{ambiguous} & \textbf{Excluded} \\
      \hline
WebQSP & 89 & 58 & 12 & 10 & 22 (5.5\%) \\ CWQ & 59 & 47 & 17 & 2 & 19 (4.8\%)
\\
      \hline
    \end{tabular}
  \end{center}
\end{table}

Two numbers must not be conflated. \Cref{tab:goldnoise} separates them: $58$ and
$47$ questions were \emph{flagged} by the consensus pass, while $22$ and $19$
were \emph{adjudicated} as defective. The gap is not noise. Most flagged
questions are all five systems converging on the same wrong answer. That
convergence is a finding about the systems rather than the labels, and it is
treated as one in \Cref{sec:echo}. A policy of rescoring on consensus alone
would have corrupted the results with exactly those cases. The excluded
questions are removed from the failure census of \Cref{sec:erroranalysis} as
benchmark artifacts rather than system failures, \textsf{ambiguous\_question} on
the same logic, an underdetermined question being no more a system failure than
a broken label. Of those excluded, $12$ of $22$ (WebQSP) and $15$ of $19$ (CWQ)
were failures of the full system. Those are the ones that shrink the census
pool.

\textbf{The accuracy tables are not rescored.} They are reported exactly as
measured, with this section as the footnote, for two reasons. Rescoring on
self-adjudicated labels would substitute one contestable ground truth for
another. And the errors cut both ways: the remaining $10$ and $4$ excluded
questions are ones the full system scored as \emph{hits} against labels now
judged broken or ambiguous. That symmetry is what makes the footnote credible
rather than defensive. It also bounds the net bias: recomputing every cell of
\Cref{tab:main} with the excluded questions removed moves each system by between
$+0.001$ and $+0.020$ Hits@1 and changes no ranking on either dataset, an order
of magnitude below the defect rates, because removing a question a system got
right subtracts from the numerator as well as the denominator.

Two scope limits belong on the record. The rule that assigned
\textsf{gold\_wrong} automatically to questions with implausibly long gold lists
is a heuristic, and its outputs were spot-checked rather than verified one by
one. And the manual adjudication covered every mixed-evidence flag and every
graph-only flag that was not a topic echo. That is broad coverage. It is not the
complete flag population.

\section{Backbone Model Selection and Qualification}\label{sec:backbone}

Every system runs on the same frozen backbone, \textsf{gpt-5.4-mini-2026-03-17},
at temperature $0.0$, with \textsf{reasoning\_effort} set explicitly to
\textsf{``none''} and the response cache enabled. Every one of the $4{,}000$
test-matrix records and every ablation record carries the model string, both
sampling parameters, and the cache mode. So any record found anywhere states
what produced it. A non-reasoning model was chosen so that sampling parameters
remain available and no hidden reasoning tokens contaminate the cost metric. The
second was verified rather than assumed, by a check that logs the
reasoning-token count on every call and read zero across every run in this
thesis.

Selection between \textsf{gpt-4.1-mini-2025-04-14} and
\textsf{gpt-5.4-mini-2026-03-17} was settled by four rounds of a determinism
probe counting how often a \emph{decision signature} --- a hash of the plan, the
expanded relations, the evaluator verdicts, and the answer --- repeats exactly
across reruns of the same questions. The rounds disagreed with one another,
because temperature-zero decoding on a hosted API is greedy rather than
deterministic and a sequential agent amplifies a single divergent token.
\Cref{app:backbone} gives the four rounds, that mechanism, and the tiebreaker
pre-specified before rounds three and four. That tiebreaker decided the choice
on longevity grounds.

\textbf{The stability figure that belongs to this thesis is the frozen model's,
which is the lower of the two.} \textsf{gpt-5.4-mini} repeated its decision
signature on $8$ of $12$ probes, or $\approx 67\%$: roughly one trajectory in
three diverging on a rerun. The $9/12$, or $\approx 75\%$, belongs to the
runner-up. At $n = 12$ neither rate is measured tightly enough to separate from
the other, and the thesis claims determinism at neither. Only the first round
survives as per-question artifacts. The qualification script opens its detail
file in truncating mode, so each later round overwrote the one before.

Two observations from those four runs matter more than the choice they settled.
Each candidate carried a \emph{persistent} ungrounded answer: 4.1-mini asserted
``United States Dollar'' for a Dominican-peso question in two of the four runs,
and 5.4-mini named the wrong Beckham child in all four. Each also carried
\emph{intermittent} ones. The same question, on the same graph, under the same
budget, moved between grounded-correct and ungrounded-wrong as the sampling
weather changed. Choosing a backbone does not remove this class of error. It
only moves it around. That is the premise of this thesis, observed over four
runs and two frontier-family models, and it is the reason
\Cref{sec:verification-layer} exists.

So reproducibility rests on the cache, not on the model.
\Cref{sec:instrumentation} specifies the mechanism, and
\Cref{app:cache-identity} gives the identity check it supports: a cache hit is
indistinguishable from a live call. A fully cached rerun therefore reproduces
the original run's budget snapshot exactly, not approximately. One defect
qualifies that. The replay restores prompt and completion tokens but not the
\textsf{reasoning\_tokens} counter. The defect does not bite here, because that
field is zero in every one of the $6{,}960$ agent runs committed for this work.
It is recorded rather than repaired, for the reason \Cref{sec:pre-registration}
gives.

Every table in \Cref{sec:results} comes from one recorded run per condition.
That run replays byte for byte for as long as the cache survives. Fresh-run
variance is a different problem. This thesis does \emph{not} measure it, the
bootstrap intervals cannot bound it, and \Cref{sec:threats} carries it as a
limitation. Latency, finally, is computed over cold-cache records only. That
restriction keeps the latency column honest where a run was interrupted and
restarted.

\section{Baseline Systems}\label{sec:baselines}

Four baselines span the space between no retrieval and agentic navigation. All
four run against the same tools, the same graph, and the same backbone.

\subsection{The Two Non-Graph Baselines}

The \textbf{no-retrieval LLM} makes a single call. It receives the question, no
facts, and an instruction to return an empty entity list when it does not know
the answer. WebQSP and CWQ predate current models and may well sit in their
training data. So this system measures the floor that every retrieval result has
to be read against. Its entity strings come from parametric memory rather than
from the graph, which means its surface forms mismatch gold more often than a
graph-grounded system's do. Its score is therefore a \emph{lower bound} on
parametric knowledge. The asymmetry is reported and not corrected, because
grading one system leniently and another strictly would corrupt the comparison.

\textbf{Vector-RAG} turns every triple into a sentence, embeds it with
MiniLM~\cite{reimers2019sentence}, and indexes it. The question retrieves the
top $30$ sentences by exact inner product, and one call answers from them. That
is dense retrieval~\cite{karpukhin2020dense} applied to a graph. Its limitation
is structural rather than incidental. A fact stored through a mediator, such as
a marriage or a film role, has no single triple that expresses it. So no single
retrieved sentence can express it either. The index drops triples with a
mediator endpoint at build time, because a verbalised half-edge means nothing
and leaks raw machine identifiers into answers. Filtering that noise out does
not repair the paradigm. Both facts are reported.

\subsection{Static GraphRAG}\label{sec:baseline-graphrag}

This is the structure-without-navigation comparison~\cite{edge2024local}. The
baseline resolves the\index{GraphRAG} question's topic entities and expands a
\emph{one-logical-hop} neighbourhood under the same relation blocklist and
mediator pass-through contract that every other graph-touching system uses. It
then reranks the result by embedding similarity to the question and answers from
the top $30$ facts in one call. It shares Vector-RAG's answering prompt and
cutoff exactly, so the difference between the two isolates retrieval strategy
and nothing else.

The radius deserves a precise statement. An earlier draft and the
implementation's own docstring both called it two hops. It is one. Expansion
runs once from each topic entity. When it meets a mediator node it hops through
and concatenates the two relations, so a mediator path costs two graph edges but
stays one logical hop. An ordinary neighbour is emitted and never expanded
again. The run records settle the point. Every logged fact reads \textsf{topic\
relation\ endpoint}. A genuine second hop would be anchored at an intermediate
entity. All $7{,}431$ logged facts begin with a topic entity.

\textbf{The radius is not the only thing that bounds this baseline. The fanout
is, and it binds harder.} Expansion retrieves at most $100$ edges per topic
entity, under a Cypher query with no \textsf{ORDER BY}. So on a higher-degree
entity the $100$ retrieved are an arbitrary sample of the neighbourhood rather
than its best or its first. Measured over the $711$ distinct entities the $800$
test questions resolve to, $55.1\%$ carry more edges than the cap retains. At
least one topic entity is truncated on $72.5\%$ of questions
(\Cref{app:fanout}). So for a hub topic the baseline answers from a small
fraction of the neighbourhood, chosen without a stated criterion. The hub-noise
diagnosis of \Cref{sec:hop-results} has to allow that the answer may never have
entered the sample the reranker saw. The full system's own
\textsf{get\_neighbors} shares the property (\Cref{app:tool-neighbors}) but not
the exposure: its cap is $200$ and applies per relation rather than per entity,
binding on $3.3\%$ of its $16{,}301$ neighbour calls.

Two implementation details fixed on development data were comparability repairs
rather than improvements. Neighbourhood expansion initially passed
\emph{through} ordinary entities and stopped at mediator nodes, exactly inverted
for Freebase and the reason machine identifiers were reaching answers. And
retrieved facts are deduplicated by endpoint pair before reranking. Cosine
similarity favours redundant phrasings of one connection and was crowding the
informative row out of the top $30$. The rule keeps the shortest rendering. That
choice resolves systematically towards the direct edge over the
mediator-expanded path whenever both connect the same pair. That bias runs
against this baseline on exactly the questions mediator expansion was added to
serve. So it is disclosed rather than left in the code.

\textbf{What this baseline licenses, and what it does not.} A one-hop retriever
cannot contain a two-hop chain. So its decay across hop strata
(\Cref{sec:hop-results}) is partly a property of the radius and not only of the
paradigm --- a distinction that has to be drawn. The alternative is to read an
implementation limit as a result. The evidence that it is a radius effect is
visible in the strata: this baseline scores $0.44$ on WebQSP's \textsf{h2},
where two-hop questions are largely mediator paths a one-hop expansion
\emph{does} reach, against $0.16$ on CWQ's \textsf{h2}, where the chains are
genuine compositions it cannot reach at all. A retriever failing uniformly on
depth would not show that gap. The claim this thesis therefore makes on the
strength of this baseline is the narrow one: \emph{retrieval fixed before
reasoning must commit to a radius in advance, and any fixed radius is wrong for
some question in the set.} It does \emph{not} claim that a two-hop static
retriever would score as this one does. That experiment was not run. The honest
expectation is that it would recover much of the CWQ \textsf{h2} gap while
losing precision to a larger and noisier window. Re-running it is the first item
in \Cref{sec:future-baselines}. Where the per-stratum argument needs a static
comparator that is definitely not radius-limited, it uses Vector-RAG, whose
one-triple context cannot represent a chain at any radius.

\subsection{Think-on-Graph (Agentic Baseline)}\label{sec:baseline-tog}

The credibility-critical comparison. Think-on-Graph~\cite{sun2024think}
is\index{Think-on-Graph} reimplemented on this environment's tools: beam width
$3$, depth $3$, language-model relation pruning, language-model entity pruning,
and a sufficiency check at each depth. It is a reimplementation rather than the
authors' code, running on this graph with this backbone. That makes it a
\emph{stronger} comparison than citing published numbers --- graph, model,
budget, and scoring rule are held constant, and the algorithm is what varies. It
correspondingly means its numbers are not comparable to the published ones and
are not presented as such.

Holding the budget constant is not the same as making it irrelevant. A fixed cap
bites harder on a method whose cost multiplies with every step. So the
comparison measures search quality and affordability together.
\Cref{sec:tog-split} separates them by splitting the results on whether the cap
actually truncated the baseline. The answer changes what the aggregate margin
should be taken to mean. Under the same $25$-call cap the full system operates
under, Think-on-Graph is clipped on $117/400$ WebQSP questions ($29.2\%$) and
$176/400$ CWQ questions ($44.0\%$). A clipped run still answers, from whatever
it gathered, by the same graceful degradation the full system uses. One
asymmetry there favours the baseline and is stated so the clip rates are not
mistaken for a count of runs that returned nothing: the baseline stops its
search one call \emph{before} the shared limit, reserving that call for the
answer it then has to compose, where AGR makes no such reservation.

\textbf{A second asymmetry is not in the baseline's favour, and it bears on how
\Cref{sec:tog-split} should be read.} Both systems call the same tools but
truncate the returned candidate sets at different widths: Think-on-Graph keeps
the first $40$ relations per entity and $20$ neighbours per relation, the
pruning widths of the reimplemented algorithm, against AGR's $300$ and $200$
(\Cref{tab:budget}). Since \textsf{get\_relations} returns rows by descending
fanout, a binding cut discards the low-fanout tail --- the material
\Cref{sec:tool-api} reports raising AGR's own cap from $80$ to $300$ to stop
discarding. It binds on $31.6\%$ of the entities Think-on-Graph expanded against
$1$ of AGR's $3{,}097$ (\Cref{app:tog-widths}). The consequence is a bound on
one specific claim. \Cref{sec:tog-split} concludes that the margin over
Think-on-Graph is affordability rather than search quality. That conclusion is
safe in the direction it is stated, since a narrower candidate set cannot
explain why the baseline runs \emph{out of calls}. But the residual gap on
unclipped questions is measured against a baseline searching a thinner pool. So
it is a lower bound on that baseline's attainable quality rather than an
estimate of it. Equalising the widths is the first item in
\Cref{sec:future-baselines}.

\subsection{Variables Held Constant, and
Certification}\label{sec:baseline-constants}

Every system shares the backbone and its sampling parameters, the knowledge
environment and its relation blocklist, the entity resolver, the question files,
the budget meter and its $25$-call cap, the output schema so that scoring reads
one field for all systems, and the record schema so that any record
self-describes with backbone, budget hash, run configuration, and source commit.
Both graph-navigating systems seed from the dataset's annotated topic entities.
That choice removes entity linking as a confound. \Cref{sec:threats} lists it as
a limitation.

Two things are \emph{not} held constant, and both are named here rather than
left buried. The first is deliberate. The baselines get a plain answering
prompt. The full system's drafting prompt carries the anti-vacuity and grounding
provisions of \Cref{sec:claim-decomposition}, and those provisions belong to the
system under test, not to the shared harness. The second is a property of the
algorithms and is where a reader should look first for a confound --- the width
of the candidate set each system considers per step, with the measured binding
rates and the bound they place on \Cref{sec:tog-split} given in
\Cref{sec:baseline-tog}, and the static baseline's own fanout cap in
\Cref{sec:baseline-graphrag}.

Baselines were debugged exclusively on development data and certified before any
test question ran. The pre-specified diagnostic targets the one failure mode
that would be an implementation defect rather than a paradigm limit ---
\emph{retrieval found the gold answer and the system hedged anyway}. That
failure mode is checkable, because both retrieval baselines log a sample of what
they retrieved. The rule, written before the numbers: under $15\%$ of hedges
certifies the baseline as-is, over $30\%$ means an over-abstinent prompt, to be
adjusted identically for both baselines with both files re-run. Measured,
Vector-RAG scored $0\%$ on both datasets and Static GraphRAG $10\%$ and $14\%$.
Both certified, with two caveats recorded rather than smoothed over. The
diagnostic samples the top ten of thirty retrieved facts. So it is a lower
bound. And Think-on-Graph's $0\%$ is \emph{vacuous} rather than clean, since it
logs no retrieved-fact sample and the diagnostic compared against an empty
string. As a multi-step navigator it was not the diagnostic's target. But the
reading is a null, not a pass. All ten test files completed at exactly $400$
records with zero failures and zero duplicate question identifiers.

\section{Metrics and the Groundedness Judge}\label{sec:metrics}

\subsection{Answer Accuracy: Hits@1 and F1}

Both metrics are computed over \emph{entity sets} after Unicode
NFKC\index{Hits"@1}\index{F1}\index{precision} normalisation, case folding, and
whitespace stripping --- no fuzzy matching, no substring containment.
\emph{Hits@1} is $1$ if any predicted entity matches any gold entity. That is
the convention in this literature, and it is what makes published numbers
comparable at all. \emph{F1} is computed per question between the predicted and
gold sets as defined in \Cref{sec:kgqa}, and averaged over questions.

The two metrics are not redundant. Hits@1 rewards a system that names one
correct entity among several wrong ones. On WebQSP that loophole is narrow on
the typical question and wide in the tail: the median question has $1.5$ gold
answers, but the mean is $7.15$ and the largest question lists $237$. The
precision component of F1 closes it. \Cref{sec:webqsp-results} shows how far
apart the two metrics can sit for a single system. That gap is a property of the
answer sets the system emits. It does not measure any component inside it:
removing the verification layer moves F1 by $+0.004$ on WebQSP and $-0.009$ on
CWQ, and neither figure is distinguishable from zero (\Cref{sec:ablations}).

Strict matching raises an obvious worry: that ``St.~Petersburg'' versus ``Saint
Petersburg'' is manufacturing failures. The worry was measured rather than
assumed. Of the full system's wrong answers, $1$ of $65$ ($2\%$) on WebQSP and
$6$ of $99$ ($6\%$) on CWQ are surface-form near-misses under aggressive
normalisation. The failure census of \Cref{sec:erroranalysis} is reading real
errors.

\subsection{Accounting for Hedges and Abstentions}\label{sec:hedge-accounting}

A \emph{hedge}\index{hedge} is an empty predicted entity set: the system
declined to assert. Hedges score zero on both accuracy metrics, but the hedge
rate gets its own column, so the trade between abstention and accuracy stays
visible. A system can buy precision with abstention. The reader has both numbers
and can watch it happen. Hedges are also not wrong answers, and
\Cref{sec:erroranalysis} keeps the two apart throughout. A wrong answer is a
reasoning error. A hedge is usually a coverage gap that never produced a
committal claim.

One counting convention follows from this, and it governs every
hits/hedges/wrong triple in the thesis: \emph{the three counts always partition
the whole set}. Nothing leaves the denominator. In particular, questions the
environment provably cannot answer (\Cref{sec:unanswerable}) count as hedges
rather than being excluded, even though the system had no route to a correct
answer. Excluding them would flatter every condition equally. It would also make
the counts incomparable with the accuracy figures, which include them.

\subsection{Two-Tier Groundedness and Hallucination
Rate}\label{sec:groundedness-metric}

Groundedness\index{groundedness} is measured at two levels. The
word\index{hallucination} \emph{tier} in this thesis refers to this metric and
nothing else.

\textbf{Tier~1, structural}, is deterministic, free, applies to every asserted
entity in all $4{,}000$ records, and operationalises
\Cref{sec:hallucination-def} directly: \emph{an asserted entity is
graph-grounded if and only if it exists as a node in the knowledge environment
and is connected to at least one of the question's topic entities within four
hops.} The entity-level ungrounded rate is the fraction of asserted entities
failing this test. The question-level rate is the fraction of answered questions
containing at least one such entity. The test is deliberately strict about
surface form --- a parametric spelling that does not exist as a node is
ungrounded \emph{by definition of the environment}. That strictness is the point
of defining hallucination relative to an environment at all.

\textbf{Tier~2, semantic}, asks the harder question Tier~1 cannot: is the entity
supported \emph{as the answer}? For a stratified sample of $60$ asserted answers
per system per dataset ($600$ judgements), the graph's own shortest connecting
paths between topic and asserted entity are retrieved. A language model judges
whether any one of them expresses the relationship the question asks about. The
governing instruction --- that \emph{mere connectedness is not support} --- is
what separates the two tiers. Two design choices are pre-specified and
consequential. Judging is against \textbf{the environment} rather than each
system's own retrieved context. The logs do not preserve per-system evidence
uniformly, and environment-judging asks every system the identical question. The
metric therefore applies even to the no-retrieval control. And path retrieval
takes the three shortest paths. That means a supporting chain can lose to
shorter irrelevant ones. Tier~2 rates are accordingly \emph{conservative lower
bounds}, uniformly depressed across all systems, and are read as a
between-system comparison rather than as absolute support rates.

\subsection{Efficiency: Tokens, LLM Calls, and Wall-Clock Latency}

Tokens and calls are counted by the shared budget meter, identically for
every\index{budget} system, and are exact. The cache replays original usage. So
they are unaffected by which runs were warm. Wall-clock latency is not, and is
computed over \textbf{cold-cache records only}. That restriction is repeated
wherever latency appears, beginning with the seconds column of
\Cref{tab:efficiency}. Where a condition's records are entirely cache replays
there is no cold-cache mean at all. That is why \Cref{tab:ablation} drops its
seconds column rather than carry one measured on four conditions of five.

\subsection{The Groundedness Judge and Its Validation}\label{sec:judge}

The Tier-2 judge sees a question, one asserted entity, and up to three
connecting paths rendered as relation chains with mediator nodes masked. It does
not see which system produced the assertion. It does not see whether the
assertion was scored a hit, what the gold answer is, or what the system's own
reasoning was. The prompt was frozen before any per-system result was inspected.
The judge runs on the same backbone as the systems it judges. \Cref{sec:threats}
states that overlap as a limitation; this chapter does not claim independence
from the model family.

One hundred judged items were drawn and written to a sheet with the judge's
verdicts stripped out. The author then labelled them \emph{blind}: neither the
judge's output file nor the answer key was opened until the sheet was saved.
Blinding is what makes the resulting statistic a validity check rather than a
measure of anchoring. The written annotation guideline fixed during calibration
is reproduced in \Cref{app:judge-rules}. \Cref{app:judge-protocol} records how
the sample was drawn and labelled. Every one of the guideline's eight rules was
forced by a case. The last of them --- treat a path expressing only a question's
identifying descriptor, rather than its actual ask, as unsupported --- is the
\emph{composition-echo} case that the failure census later finds at scale
(\Cref{sec:echo}).

On $n = 100$ items, observed agreement is $85\%$, with near-symmetric
marginals\index{Cohen's kappa} ($54$ human-supported against $51$
judge-supported) and near-symmetric disagreement ($9$ human-yes/judge-no against
$6$ human-no/judge-yes). So neither rater is systematically the more lenient.
Cohen's $\kappa$~\cite{cohen1960coefficient} is \[ \kappa = \frac{0.8500 -
0.5008}{1 - 0.5008} = 0.6995 . \]

The pre-specified threshold was $\kappa \geq 0.7$. Rounded to three decimals
this reads $0.700$ and appears to clear. It does not: it falls short by
$0.0005$, and it is reported as short. What follows is a matter of weight rather
than of discarding. By the conventional bands~\cite{landis1977measurement} a
$\kappa$ of $0.70$ is substantial agreement. The task is genuinely hard: whether
a relation chain really expresses an asked relationship has real grey area for a
careful human, as the guideline's list of decided edge cases attests. The Tier-2
numbers are therefore reported with the $\kappa$ beside them and the miss
stated. No conclusion rests on Tier~2 alone --- where it is load-bearing, in the
between-system semantic-support comparison of \Cref{sec:groundedness-results},
the finding is also supported by the precision column of the main matrix, which
is measured without a judge. Tightening the judge prompt and re-judging a fresh
sample until the bar was cleared was available and pre-committed. It was not
taken, because the honest number was already in hand, and re-running until a
pre-specified threshold is met is exactly what pre-specification exists to
prevent.

\section{Ablation Conditions and the Development
Set}\label{sec:ablation-conditions}

Four conditions, each a single deviation from the frozen configuration, each
isolating one mechanism:

\begin{itemize}
  \item \textbf{no-planner} --- decomposition disabled; the whole question
becomes one objective. Anchor seeding is untouched, so this varies decomposition
and nothing else.
  \item \textbf{no-backtrack} --- \textsf{max\_backtracks} $= 0$; a
configuration change, which correctly alters the budget hash stamped on the
records.
  \item \textbf{no-verifier} --- \textsf{verify\_claims} $=$ false; the drafting
call is unchanged, claims included, but no claim is checked and the draft is
emitted verbatim (\Cref{sec:repair}), so the condition varies the checking and
not the drafting.
  \item \textbf{embedding-only} --- $\alpha = 1.0$, removing the language-model
term from the hybrid scorer while leaving $\tau$ at its frozen value.
\end{itemize}

The reference row is the \emph{unmodified} system restricted to the same
questions, obtained by filtering the main-matrix records to the half-set
question identifiers rather than by re-running them. That filtering also
supplies the matched pairs McNemar's test requires. Comparing against the full
$400$-question score would be invalid, accuracy on $400$ questions and on $198$
being different numbers for the same system. Ablations run on stratified halves
--- $200$ WebQSP and $198$ CWQ --- as the pre-specified response to the
benchmark's burn coming in above projection, cutting the ablation \emph{sample}
rather than any system or condition. The odd $198$ is floor division within each
stratum, CWQ's $137/211/49/3$ halving to $68/105/24/1$. The consequence, reduced
statistical power, is confronted in \Cref{sec:ablations}.

\subsection{The Development Set and the
\texorpdfstring{$\alpha$--$\tau$}{alpha-tau} Sweep}\label{sec:hyperparams}

\label{sec:dev-construction}%
The development set is $80$ questions, disjoint from both test samples by an
assertion in the test-set builder. Seventy-seven were mined from the train and
validation splits to span hop counts and to include mediator-heavy and
constraint-bearing questions, and certified for answer reachability by the
procedure applied to the test samples. The other three are hand-written and,
having no gold answer to reach, fall outside both. That is why the coverage
report carries $77$ rows. The set \emph{contains} the twenty-question smoke set
of \Cref{sec:design-validation}. That is why identifiers such as
\textsf{smoke-20} appear in the dev-set listing of \Cref{app:samples}. Seven
questions are unanswerable, in two buckets that are not interchangeable: three
hand-written ones naming a fabricated entity, so that a correct refusal can be
distinguished from a coverage failure, and four ordinary benchmark questions
with real gold answers this environment cannot express
(\Cref{sec:unanswerable}).

The two free parameters of the navigation loop are the blend weight $\alpha$ in
the hybrid scorer (\Cref{sec:scoring}) and the low-signal backtracking threshold
$\tau$ (\Cref{sec:backtracking}). They were swept jointly over $\alpha \in
\{0.3, 0.5, 0.7, 1.0\}$ and $\tau \in \{0.0, 0.20\}$ --- eight conditions, each
a complete run of all $80$ questions, with two additional draft-only conditions
to price the verification layer. Each condition is a constructor argument whose
values are stamped into every record it writes. So no condition can be
mislabelled after the fact. The chosen values, $\alpha = 0.7$ and $\tau = 0.20$,
were frozen on the basis of that sweep and never revisited.
\Cref{sec:dev-tuning} reports it.

\subsection{Implementation, Environment, and
Reproducibility}\label{sec:implementation}

The system is implemented in Python against LangGraph for the state machine, the
Neo4j Python driver for graph access~\cite{robinson2015graph}, and
sentence-transformers for embeddings~\cite{reimers2019sentence}, with the graph
in a local Neo4j instance. All Cypher~\cite{francis2018cypher} is templated
inside the tool layer and never composed by a language model
(\Cref{sec:tool-api}). Randomised operations use seed $42$ throughout.

Three practices carry more reproducibility weight than the dependency list.
\Cref{app:implementation} records them as techniques. \emph{Records
self-describe}: every record carries the backbone description, a
budget-configuration hash, the full run configuration, and the source commit.
That self-description is what made the cache-replay incident of
\Cref{sec:backbone} diagnosable. \emph{Frozen artifacts are committed and
regenerable ones are not}, the question sets and result files embodying spent
money and non-repeatable runs where the graph import and the indexes are
deterministic products of committed scripts. And \emph{code changes are proved
harmless by cache identity} rather than by a passing test suite: re-running an
answered question and confirming every call is a cache hit proves byte-identical
prompts and hence an untouched frozen path. That proof is the certificate behind
the ablation reference row.

Two limits belong here rather than in a footnote. The archived results carry
neither per-folder documentation nor generated manifests. So file provenance
rests on the self-describing records and on this chapter. And one generated
summary in the development-set archive predates a fix applied to the runs it
summarises. So rescoring that directory yields different numbers than the
committed summary. The discrepancy stands. \Cref{sec:dev-tuning} states it where
those numbers are used, because editing a generated file is not a fix.

\section{Main Results}\label{sec:results}

This part reports the measurements that \Cref{sec:setup} specified, arranged so
that the weakest claims sit next to the strongest. The headline accuracy result
is large and significant on every comparison. The structural groundedness result
is emphatic but points to a different cause than the one expected. The semantic
groundedness result is underpowered. Three of the four ablations detect nothing.
All four outcomes get the same volume.

Every test-set number comes from \textsf{results/phase4/thesis\_numbers.json}.
The development sweep of \Cref{sec:dev-tuning} is the one set of figures that
does not: it reads the archived phase-3 runs directly, as \Cref{ch:framework}
does. Confidence intervals are $95\%$ bootstrap intervals over questions, from
$10{,}000$ resamples at seed~$42$~\cite{efron1979bootstrap}. Paired comparisons
use the exact-binomial form of McNemar's test~\cite{mcnemar1947note} on
per-question correctness.

\subsection{Development-Set Tuning Outcomes}\label{sec:dev-tuning}

The sweep varied the blend weight $\alpha$ and the low-signal backtracking
threshold $\tau$ jointly over the $80$-question development set
(\Cref{sec:hyperparams}). \Cref{tab:sweep} reports all eight conditions, plus
the two draft-only conditions that put a price on the verification layer.

\begin{table}[!tb]
  \begin{center}
    \caption[The development-set sweep]{The development-set sweep, $n = 80$.
    Hits, hedges, and wrong answers add up to the whole set. \textsf{Verify}
    means claim verification was enabled. The two draft-only rows disable it.
    All rows come from the sweep as it was run, before the correction to the
    answer-entity filter of \Cref{sec:entity-filtering}. The text gives the
    frozen condition's post-correction figures.}
    \label{tab:sweep}
    \small
    \begin{tabular}{|r|r|c|r|r|r|r|r|r|r|}
      \hline
      $\boldsymbol{\alpha}$ & $\boldsymbol{\tau}$ & \textbf{Verify} &
      \textbf{Hits} & \textbf{Hedge} & \textbf{Wrong} & \textbf{F1} &
      \textbf{Tokens} & \textbf{Calls} & \textbf{Bktrk} \\
      \hline
      0.3 & 0.00 & yes & 61 & 12 & 7 & 0.728 & 4{,}999 & 7.4 & 16 \\
      0.3 & 0.20 & yes & 61 & 12 & 7 & 0.728 & 5{,}030 & 7.5 & 18 \\
      0.5 & 0.00 & yes & 62 & 10 & 8 & 0.743 & 4{,}874 & 7.2 & 16 \\
      0.5 & 0.20 & yes & 62 & 10 & 8 & 0.743 & 4{,}874 & 7.2 & 16 \\
      \textbf{0.7} & 0.00 & yes & 64 & 11 & 5 & 0.769 & 4{,}925 & 7.3 & 17 \\
      \textbf{0.7} & \textbf{0.20} & \textbf{yes} & \textbf{64} & 11 & 5 & 0.769
        & 4{,}925 & 7.3 & 17 \\
      1.0 & 0.00 & yes & 59 & 15 & 6 & 0.713 & 3{,}408 & 4.9 & 22 \\
      1.0 & 0.20 & yes & 59 & 16 & 5 & 0.713 & 3{,}400 & 4.9 & 26 \\
      \hline
      0.5 & 0.20 & no  & 64 &  9 & 7 & 0.770 & 4{,}807 & 7.1 & 16 \\
      0.7 & 0.20 & no  & 66 & 10 & 4 & 0.797 & 4{,}858 & 7.2 & 17 \\
      \hline
    \end{tabular}
  \end{center}
\end{table}

The sweep supports three readings. Only the first is the one it was run to
obtain: $\alpha = 0.7$ is the maximum. \Cref{app:sweep} gives the other two. The
curve is not clean. The wrong-answer column peaks at $\alpha = 0.5$, where the
system commits more often, and $\tau$ changes almost nothing at the values
swept. That flatness is what the signal-maximum formulation of
\Cref{sec:backtracking} predicts. Both facts are reported rather than smoothed
away. In particular, $\tau = 0.20$ was frozen on the strength of a sweep in
which it barely mattered. That is stated plainly rather than dressed up as a
tuned choice.

\textbf{Verification costs accuracy on this set.} Against the frozen condition,
the draft-only row scores $66$ hits to $64$ and F1 $0.797$ to $0.769$. After the
answer-entity filter of \Cref{sec:entity-filtering} was corrected --- a change
made on development data before any test question ran --- the frozen condition
scores $65$ hits, $11$ hedges, $4$ wrong, F1 $0.785$. Across those two rows the
layer costs one correct answer and removes no wrong ones on this set of $80$.
That comparison spans the correction, since the draft-only condition was never
re-run after it. It is therefore the closest available reading rather than a
matched pair. \Cref{sec:ablations} finds the same thing at $n \approx 200$ on
test data. \Cref{sec:findings} explains why that is the expected result rather
than a disappointing one. Only the frozen condition was re-run after the
correction. So \Cref{tab:sweep} is the pre-correction sweep throughout ---
internally consistent, since all eight conditions ran under the same code. Where
the committed summary and the archived per-question records disagree, the
records are authoritative.

\subsection{Main Results on WebQSP}\label{sec:webqsp-results}

\begin{table}[!tb]
  \begin{center}
    \caption[Main results on both test sets]{Main results, $n = 400$ per
    dataset. Brackets are $95\%$ bootstrap intervals. \textbf{P} and \textbf{R}
    are mean per-question precision and recall. \textbf{Hedge} is the share of
    questions on which the system asserted nothing. The Hits@1 ceilings are
    $97.0\%$ and $99.2\%$ (\Cref{tab:testsets}). Every figure here is measured
    inside the environment of \Cref{ch:environment}. That environment's
    per-question construction is what puts the ceilings so high. So the levels
    are not comparable with published results over full Freebase.
    \Cref{sec:threats} states why.}
    \label{tab:main}
    \small
    % Seven columns of bracketed intervals leave this a point or so wider than
    % the text block; trimming the intercolumn padding is enough to seat it.
    \setlength{\tabcolsep}{5pt}
    \begin{tabular}{|l|l|l|l|r|r|r|}
      \hline
      \textbf{Data} & \textbf{System} & \textbf{Hits@1} & \textbf{F1} &
      \textbf{P} & \textbf{R} & \textbf{Hedge} \\
      \hline
      WebQSP & No-retrieval  & 0.453 [0.403, 0.500] & 0.305 [0.266, 0.345] &
        0.380 & 0.303 & 12.2\% \\
      WebQSP & Vector-RAG    & 0.568 [0.517, 0.618] & 0.432 [0.389, 0.475] &
        0.469 & 0.460 & 27.0\% \\
      WebQSP & GraphRAG      & 0.340 [0.292, 0.388] & 0.262 [0.224, 0.302] &
        0.279 & 0.284 & 55.8\% \\
      WebQSP & Think-on-Graph & 0.660 [0.613, 0.705] & 0.477 [0.436, 0.518] &
        0.571 & 0.480 & 18.8\% \\
      WebQSP & \textbf{AGR}  & \textbf{0.755} [0.713, 0.797] & \textbf{0.642}
        [0.600, 0.684] & \textbf{0.697} & \textbf{0.653} & \textbf{8.2\%} \\
      \hline
      CWQ & No-retrieval  & 0.307 [0.263, 0.352] & 0.279 [0.236, 0.323] & 0.285
        & 0.282 & 38.0\% \\
      CWQ & Vector-RAG    & 0.203 [0.163, 0.242] & 0.168 [0.133, 0.204] & 0.171
        & 0.189 & 48.2\% \\
      CWQ & GraphRAG      & 0.205 [0.168, 0.245] & 0.183 [0.147, 0.222] & 0.187
        & 0.194 & 60.8\% \\
      CWQ & Think-on-Graph & 0.435 [0.388, 0.482] & 0.378 [0.332, 0.423] & 0.403
        & 0.384 & 22.5\% \\
      CWQ & \textbf{AGR}  & \textbf{0.522} [0.475, 0.570] & \textbf{0.469}
        [0.423, 0.514] & \textbf{0.484} & \textbf{0.481} & 23.0\% \\
      \hline
    \end{tabular}
  \end{center}
\end{table}

On WebQSP, AGR answers $302$ of $400$ questions correctly. That is $75.5\%$
against a\index{Hits"@1}\index{WebQSP} reachability ceiling of $97.0\%$, $+9.5$
points over Think-on-Graph, and $+30.2$ over the parametric floor. The floor
matters. WebQSP is a decade old and quite possibly in the backbone's pretraining
data. The no-retrieval control confirms as much: the backbone answers $45.3\%$
of these questions from memory alone. Any claim about what retrieval contributes
therefore starts from $0.453$, not from zero.

\textbf{F1 widens the margin that Hits@1 understates.} Against Think-on-Graph
the Hits@1 gap is $+9.5$ points; the F1 gap is $+16.5$ ($0.642$ against
$0.477$). The precision column shows where the extra margin comes from: $0.697$
against $0.571$. WebQSP questions carry seven gold answers on average, so a
system that asserts a long list collects Hits@1 credit for the one entity that
matches and pays nothing for the rest. Precision charges for the rest. Together
the two columns support the sentence this thesis leads with: AGR answers more
questions and asserts fewer falsehoods. It does so at an $8.2\%$ hedge rate, the
lowest of the five systems, so the precision is not bought with abstention.

\textbf{Raw hits mislead on the retrieval baselines. The correction is a
finding.} GraphRAG scores $0.340$, below the no-retrieval control's $0.453$.
That gap invites the conclusion that the implementation is broken. Put the
wrong-answer counts beside the hits and the picture inverts. No-retrieval
asserts on $351$ of $400$ questions and is wrong on $170$ of them --- $51.6\%$
\emph{assertion precision}\index{assertion precision}, the share of the
questions a system commits on that it answers correctly. It is a per-question
measure. It is also distinct from the per-entity precision of \Cref{tab:main},
which is averaged over all $400$ questions whether the system committed on them
or not. GraphRAG asserts on $177$ and is wrong on $41$ --- $76.8\%$. The
mechanism is the prompt: the retrieval baselines are instructed to answer only
from the facts given and to return nothing otherwise. When retrieval misses,
they therefore abstain, while the no-retrieval control has no such constraint
and guesses from memory. On a memorised benchmark, guessing pays. \emph{A
grounded system scoring below an ungrounded one on raw hits while committing
four times fewer falsehoods} is the contamination story the control was built to
expose. It is reported as a result rather than as an anomaly.

\subsection{Main Results on ComplexWebQuestions}\label{sec:cwq-results}

On CWQ, AGR answers $209$ of $400$ questions: $52.2\%$ against a $99.2\%$
ceiling, $+8.7$\index{ComplexWebQuestions}\index{precision} points over
Think-on-Graph and $+21.5$ over the floor. Every system scores lower here than
on WebQSP. But the ordering changes in one important way: \emph{both} static
retrieval baselines now fall below the parametric control. Vector-RAG drops from
$0.568$ to $0.203$ and GraphRAG from $0.340$ to $0.205$, while no-retrieval
falls only from $0.453$ to $0.307$.

The reason is structural rather than incidental, though the two baselines earn
that description to different degrees. Vector-RAG's context is a single
verbalised triple. One triple cannot hold a chain, at any retrieval quality or
any radius. \Cref{sec:baselines} predicted that paradigm limit at implementation
time, and \Cref{sec:hop-results} shows it stratum by stratum. GraphRAG's context
is a one-logical-hop neighbourhood (\Cref{sec:baseline-graphrag}). Its CWQ
number is the weaker evidence of the two, because it confounds the paradigm with
the radius: part of its fall belongs to a boundary chosen in advance, not to
static retrieval as such. The claim therefore rests on Vector-RAG.

The number worth leading with on this dataset is assertion precision. AGR and
Think-on-Graph commit at effectively the same rate: $308$ and $310$ of $400$
questions answered rather than hedged. From those near-identical commitments,
AGR is right $209$ times and wrong $99$, against $174$ and $136$ for
Think-on-Graph. That is $67.9\%$ assertion precision against $56.1\%$. On the
harder, deeper benchmark, the system with an explicit plan, a guided scorer, and
a claim check finds $35$ more answers and asserts $37$ fewer falsehoods than
beam search does on the same graph, with the same model, from the same number of
commitments. \Cref{sec:tog-split} attributes that margin to the budget rather
than to resolution quality.

\begin{table}[!tb]
  \begin{center}
    \caption[Paired McNemar tests against each baseline]{Paired McNemar tests of
    AGR against each baseline, on per-question correctness over the same $400$
    questions. \textbf{AGR only} counts the questions that AGR alone answers
    correctly; \textbf{Baseline only} counts the reverse. Every comparison
    rejects.}
    \label{tab:mcnemar}
    \small
    \begin{tabular}{|l|l|r|r|l|}
      \hline
      \textbf{Dataset} & \textbf{Baseline} & \textbf{AGR only} &
      \textbf{Baseline only} & \textbf{$p$} \\
      \hline
      WebQSP & No-retrieval   & 152 & 31 & $2.3 \times 10^{-20}$ \\
      WebQSP & Vector-RAG     & 117 & 42 & $2.2 \times 10^{-9}$ \\
      WebQSP & GraphRAG       & 186 & 20 & $6.5 \times 10^{-35}$ \\
      WebQSP & Think-on-Graph &  76 & 38 & $4.8 \times 10^{-4}$ \\
      \hline
      CWQ & No-retrieval   & 121 & 35 & $2.7 \times 10^{-12}$ \\
      CWQ & Vector-RAG     & 151 & 23 & $2.8 \times 10^{-24}$ \\
      CWQ & GraphRAG       & 148 & 21 & $1.0 \times 10^{-24}$ \\
      CWQ & Think-on-Graph &  86 & 51 & $3.5 \times 10^{-3}$ \\
      \hline
    \end{tabular}
  \end{center}
\end{table}

\Cref{tab:mcnemar} settles the comparisons that matter. All eight reject. That
includes the two against the agentic baseline, the only comparison in the table
where both systems navigate the same graph, with the same model, under the same
call budget. The discordant counts make the margin concrete rather than
distributional: on CWQ, AGR uniquely solves $86$ questions Think-on-Graph misses
against $51$ the other way. What that budget does to the comparison is the
subject of the next section. It changes what the margin means.

\section{Where the Margin Over the Agentic Baseline Comes
From}\label{sec:tog-split}

The shared budget is a control. Like any control, it is doing something.
The\index{Think-on-Graph} call cap binds on Think-on-Graph far more often than
on AGR: beam search costs width times depth. \Cref{sec:efficiency} records the
clip rates. So ``same budget'' does not mean ``same freedom to search''.
Splitting the $400$ questions of each dataset on whether the cap actually cut
Think-on-Graph off separates the two things the aggregate confounds: how good
its search is, and whether it can afford to run that search to completion.

\begin{table}[!tb]
  \begin{center}
    \caption[The agentic-baseline comparison, split on budget clipping]{The
    comparison against the agentic baseline, split on whether the shared
    $25$-call cap cut it short. Clipping is read from the per-record
    \textsf{budget\_exhausted} trace flag, not from a call count, because a run
    may spend its last permitted call on the step that finishes it. The
    \emph{combined} rows reproduce the Hits@1 column of \Cref{tab:main},
    recomputed independently from the raw records. Generated from
    \textsf{results/phase4/thesis\_numbers.json}.}
    \label{tab:tog-split}
    \small
    % GENERATED by scripts/build_figures.py from
% results/phase4/thesis_numbers.json -- do not edit.

\begin{tabular}{|l|l|r|r|r|r|}
\hline
\textbf{Dataset} & \textbf{Subset} & \textbf{$n$} & \textbf{ToG} & \textbf{AGR} & \textbf{$\Delta$} \\
\hline
WebQSP & ToG ran to completion & 283 & 0.852 & 0.788 & $-0.064$ \\
 & ToG clipped by the cap & 117 & 0.197 & 0.675 & $+0.479$ \\
& \emph{combined} & 400 & \emph{0.660} & \emph{0.755} & $\mathbf{+0.095}$ \\
\hline
CWQ & ToG ran to completion & 224 & 0.629 & 0.607 & $-0.022$ \\
 & ToG clipped by the cap & 176 & 0.188 & 0.415 & $+0.227$ \\
& \emph{combined} & 400 & \emph{0.435} & \emph{0.522} & $\mathbf{+0.087}$ \\
\hline
\end{tabular}

  \end{center}
\end{table}

\textbf{On the questions it is allowed to finish, Think-on-Graph is ahead on
both datasets.} It scores $0.852$ against AGR's $0.788$ on WebQSP and $0.629$
against $0.607$ on CWQ. The entire aggregate margin comes from the $29.2\%$ and
$44.0\%$ of questions where the cap truncates it. There its accuracy falls to
$0.197$ and $0.188$ while AGR's holds at $0.675$ and $0.415$. This is stated
plainly for two reasons. It is the first thing an unsympathetic reading of
\Cref{tab:mcnemar} should look for. And burying it would misrepresent what was
measured.

This does not weaken the result. It changes what the result says. The finding is
not that guided best-first search resolves entities better than
language-model-pruned beam search. On the evidence here it does not, and the
unclipped subset is the fairest available estimate of that. The finding is that
\emph{beam search cannot afford to finish under a budget a deployed system would
actually impose}, and \emph{guided search can}. Width times depth multiplies,
and every question pays it. A plan plus a scored frontier does not multiply. A
comparison run without a cap would measure search quality in the abstract. That
would be a legitimate experiment, and a different one. This thesis fixed the
budget in advance (\Cref{sec:baselines}) because cost is a property of the
method, not a nuisance to control away. On that reading the clip rate is not an
artefact of the setup. It is the result. \Cref{sec:efficiency} prices the same
finding in tokens.

Three qualifications keep the claim honest. The first is that the cap of $25$
calls is itself a choice. A larger one would move the split --- the clip rates
are reported so a reader can see how far. The reimplementation's beam width and
depth follow the original paper. So its cost profile is the published
algorithm's rather than a weakened variant. That fidelity is what makes the cost
finding meaningful instead of circular. But the candidate widths \emph{beneath}
the beam do not come from the paper: the baseline prunes from $40$ relations and
$20$ neighbours where AGR prunes from $300$ and $200$. Those cuts bind on about
a third of its expansions against almost none of AGR's
(\Cref{sec:baseline-tog}). That asymmetry cannot explain the clipping --- a
thinner candidate set makes a step cheaper, not dearer. So the affordability
finding stands. It does bear on the sentence before it: $0.852$ and $0.629$ on
the unclipped subset are a \emph{lower bound} on what this baseline's search
could resolve given AGR's candidate widths, not an estimate of it. The claim
that guided search does not out-resolve beam search is correspondingly the
weaker one.

\textbf{The breadth of AGR's answers is checked rather than assumed.} On
answered WebQSP questions AGR names $3.37$ entities against Think-on-Graph's
$2.56$. A metric scoring a hit on any match rewards naming more. The check is
entity precision over those same answered questions --- \Cref{tab:main}'s
measure restricted to the answered subset, not the per-question assertion
precision of \Cref{sec:cwq-results} --- where AGR leads $0.760$ to $0.703$. On
CWQ the figures are $1.54$ against $1.40$ entities and $0.629$ against $0.520$
precision. If AGR bought breadth without accuracy, precision would fall as the
entity count rose. It does not fall.

\section{Positioning Against Published RoG Results}\label{sec:rog-comparison}

Two facts in this thesis invite a direct comparison. The knowledge environment
of \Cref{sec:source-data} is built from the per-question subgraphs\index{RoG}
that the RoG distribution publishes, and \Cref{sec:path-kgqa} surveys
RoG~\cite{luo2024reasoning} as the closest published system in the
path-retrieval family. A reader who notices both will ask how the two systems
compare. The aggregate answer does not favour this work.

\begin{table}[!tb]
  \begin{center}
    \caption[AGR against RoG's published figures]{AGR against RoG's published
    figures, in points. RoG's row is Table~1 of \cite{luo2024reasoning},
    measured over the full test splits ($1{,}628$ and $3{,}531$ questions).
    AGR's row is \Cref{tab:main}, measured over the certified $400$-question
    samples of \Cref{tab:testsets}. Those samples are stratified draws from
    those same splits. So the two rows are computed over nested populations, but
    under different protocols. The paragraphs below say what that permits and
    what it forbids.}
    \label{tab:rog}
    \small
    \begin{tabular}{|l|c|c|c|c|}
      \hline
      & \multicolumn{2}{c|}{\textbf{WebQSP}} & \multicolumn{2}{c|}{\textbf{CWQ}}
      \\
      \cline{2-5}
      \textbf{System} & \textbf{Hits@1} & \textbf{F1} & \textbf{Hits@1} &
                      \textbf{F1} \\
      \hline
      RoG (published)  & $85.7$ & $70.8$ & $62.6$ & $56.2$ \\
      \hline
      AGR (this work)  & $75.5$ & $64.2$ & $52.2$ & $46.9$ \\
      \hline
    \end{tabular}
  \end{center}
\end{table}

\textbf{RoG leads on all four figures. The margin is not sampling noise.} AGR's
$95\%$ confidence intervals on Hits@1 are $[71.3, 79.7]$ on WebQSP and $[47.5,
57.0]$ on CWQ. RoG's $85.7$ and $62.6$ fall outside both. That is stated before
the qualifications. The qualifications do not reverse it, and a comparison
offered only with its excuses attached is not a comparison.

Three differences bound what \Cref{tab:rog} establishes. The first is decisive.
The other two are ordinary.

\textbf{RoG is trained on these benchmarks. AGR is trained on nothing.}
Section~5.1 of \cite{luo2024reasoning} states that its backbone is ``instruction
finetuned on the training split of WebQSP and CWQ as well as Freebase for 3
epochs''. In the distribution this thesis also uses, those training splits hold
$2{,}826$ and $27{,}639$ questions. AGR performs no training of any kind: the
backbone is frozen, prompted at temperature $0$ (\Cref{sec:backbone}), and the
system has seen no question from either benchmark before it is evaluated. The
comparison is therefore between a supervised system and a zero-shot one. That is
a difference in kind rather than in degree. It is the same objection
\Cref{tab:reported} already raises against every fine-tuned entry in the
surveyed literature.

\textbf{The backbones differ.} RoG fine-tunes LLaMA2-Chat-7B. AGR prompts the
model named in \Cref{sec:backbone}. \Cref{sec:comparison-table} argues that
backbone accounts for much of the spread in cross-paper tables. That is
precisely why \Cref{sec:setup} holds it constant across the five systems it
measures. It cannot be held constant against a system whose contribution
\emph{is} its fine-tuned weights. So no protocol available here would make this
row comparable in the sense \Cref{tab:main} uses the word.

\textbf{The evaluation populations differ.} RoG reports over the full test
splits. The figures here are over stratified $400$-question samples of the same
splits (\Cref{sec:test-sets}). The samples preserve the hop-stratum proportions
of their populations. So they estimate the same quantity. The intervals quoted
above are the estimate's uncertainty. RoG publishes no interval. So the gap can
be bounded from one side only.

What the comparison does settle is worth stating positively. It fixes AGR's
place in the published landscape: a training-free system reaching $75.5$ and
$52.2$ against a supervised system's $85.7$ and $62.6$, on the same data, with
no in-domain examples. And it locates this thesis's contribution somewhere other
than the accuracy column. RoG's answers are read off retrieved paths. That makes
them faithful by construction. The pre-generation verification question of
\Cref{sec:verification-layer} never comes up for them (\Cref{sec:path-kgqa}).
The price is that the plan is generated once, before retrieval, with no
mechanism to revise it when it fails to instantiate. AGR's claims are the
opposite shape: generated freely and then checked against the graph the agent
actually walked. That check is what makes an unsupported assertion detectable at
all. Neither the component-level ablation of \Cref{sec:ablations}, with the
stratum-dependent decomposition finding it produced, nor the failure census of
\Cref{sec:erroranalysis} has a counterpart in RoG's evaluation. None of them
would be settled by closing the gap in \Cref{tab:rog}.

\section{Accuracy Broken Down by Hop Count}\label{sec:hop-results}

\begin{table}[!tb]
  \begin{center}
    \caption[Hits@1 and F1 per hop stratum]{Hits@1 and F1 per hop stratum.
    Stratum sizes are those of \Cref{tab:testsets}. WebQSP's \textsf{h3plus} has
    $n = 4$ and is not interpreted.}
    \label{tab:strata}
    \small
    \begin{tabular}{|l|l|l|l|l|l|}
      \hline
      \textbf{Data} & \textbf{System} & \textbf{h1} & \textbf{h2} &
      \textbf{h3plus} & \textbf{unreach.} \\
      \hline
      WebQSP & No-retrieval   & 0.43 / 0.29 & 0.52 / 0.36 & 0.25 / 0.00 & 0.17 /
        0.17 \\
      WebQSP & Vector-RAG     & 0.82 / 0.63 & 0.12 / 0.07 & 0.00 / 0.00 & 0.08 /
        0.08 \\
      WebQSP & GraphRAG       & 0.30 / 0.23 & 0.44 / 0.35 & 0.25 / 0.07 & 0.08 /
        0.08 \\
      WebQSP & Think-on-Graph & 0.63 / 0.45 & 0.78 / 0.58 & 0.25 / 0.07 & 0.17 /
        0.17 \\
      WebQSP & \textbf{AGR}   & 0.75 / 0.64 & \textbf{0.84 / 0.73} & 0.00 / 0.00
        & 0.08 / 0.08 \\
      \hline
      CWQ & No-retrieval   & 0.38 / 0.35 & 0.30 / 0.27 & 0.14 / 0.14 & 0.00 /
        0.00 \\
      CWQ & Vector-RAG     & 0.33 / 0.27 & 0.14 / 0.11 & 0.10 / 0.09 & 0.67 /
        0.67 \\
      CWQ & GraphRAG       & 0.34 / 0.30 & 0.16 / 0.14 & 0.02 / 0.02 & 0.33 /
        0.33 \\
      CWQ & Think-on-Graph & \textbf{0.53} / 0.48 & 0.37 / 0.32 & 0.45 / 0.32 &
        0.33 / 0.33 \\
      CWQ & \textbf{AGR}   & 0.46 / 0.42 & \textbf{0.55 / 0.49} & \textbf{0.57 /
        0.50} & 0.33 / 0.33 \\
      \hline
    \end{tabular}
  \end{center}
\end{table}

\begin{figure}[!tb]
  \begin{center}
    % GENERATED by scripts/build_figures.py from
% results/phase4/thesis_numbers.json -- do not edit.

\begin{tikzpicture}
\begin{axis}[agrplot, name=hop1, 
  title={WebQSP},
  xlabel={Hop stratum}, ylabel={Hits@1},
  xtick={0,1,2}, xticklabels={1 hop ($n{=}256$), 2 hops ($n{=}128$), 3+ hops ($n{=}4$)},
  xmin=-0.25, xmax=2.25, ymin=0, ymax=0.9,
  legend style={at={(1.08,-0.30)}, anchor=north, legend columns=5, column sep=1.8mm},
  legend cell align=left,
  legend image post style={solid},
]
\addplot[dashed, mark=o, mark size=2.2pt, color=black!55, thick] coordinates {(0,0.43) (1,0.52) (2,0.25)};
\addlegendentry{No-retrieval}
\addplot[dashed, mark=square, mark size=2.2pt, color=black!70, thick] coordinates {(0,0.82) (1,0.12) (2,0.0)};
\addlegendentry{Vector-RAG}
\addplot[dashed, mark=triangle, mark size=2.2pt, color=black!70, thick] coordinates {(0,0.3) (1,0.44) (2,0.25)};
\addlegendentry{GraphRAG}
\addplot[dashed, mark=diamond, mark size=2.2pt, color=black!85, thick] coordinates {(0,0.63) (1,0.78) (2,0.25)};
\addlegendentry{Think-on-Graph}
\addplot[dashed, mark=*, mark size=2.2pt, color=black, thick] coordinates {(0,0.75) (1,0.84) (2,0.0)};
\addlegendentry{AGR}
\end{axis}
\begin{axis}[agrplot, name=hop2, at={($(hop1.east)+(12mm,0)$)}, anchor=west,
  title={ComplexWebQuestions},
  xlabel={Hop stratum}, ylabel={Hits@1},
  xtick={0,1,2}, xticklabels={1 hop ($n{=}137$), 2 hops ($n{=}211$), 3+ hops ($n{=}49$)},
  xmin=-0.25, xmax=2.25, ymin=0, ymax=0.9,
%
]
\addplot[mark=o, mark size=2.2pt, color=black!55, thick] coordinates {(0,0.38) (1,0.3) (2,0.14)};
\addplot[mark=square, mark size=2.2pt, color=black!70, thick] coordinates {(0,0.33) (1,0.14) (2,0.1)};
\addplot[mark=triangle, mark size=2.2pt, color=black!70, thick] coordinates {(0,0.34) (1,0.16) (2,0.02)};
\addplot[mark=diamond, mark size=2.2pt, color=black!85, thick] coordinates {(0,0.53) (1,0.37) (2,0.45)};
\addplot[mark=*, mark size=2.2pt, color=black, thick] coordinates {(0,0.46) (1,0.55) (2,0.57)};
\end{axis}
\end{tikzpicture}

  \end{center}
  \caption[Hits@1 across hop strata]{Hits@1 across hop strata. The shape is the
  finding: on ComplexWebQuestions AGR is the only system whose accuracy rises
  with hop count. Every other profile ends below where it started --- three of
  them monotonically, Think-on-Graph by falling and partly recovering. WebQSP is
  drawn dashed because its \textsf{h3plus} stratum holds four questions and
  cannot carry a trend. The stratum sizes are printed on the axis for that
  reason. The \textsf{unreachable} stratum is omitted here and discussed below,
  since by construction it measures assertion rather than accuracy. Generated
  from \textsf{results/phase4/thesis\_numbers.json}.}
  \label{fig:hop-strata}
\end{figure}

\Cref{tab:strata} is the direct answer to RQ1. It answers it with a shape rather
than with a difference of means --- \Cref{fig:hop-strata} is that
shape.\index{hop stratum}

\textbf{On CWQ, AGR's accuracy rises with hop count.} Hits@1 goes $0.46 \to 0.55
\to 0.57$ across \textsf{h1}, \textsf{h2}, \textsf{h3plus}. F1 goes $0.42 \to
0.49 \to 0.50$. It is the only system on that dataset that ends above where it
started. The other four all end below. Three of them decay monotonically:
No-retrieval $0.38 \to 0.30 \to 0.14$, Vector-RAG $0.33 \to 0.14 \to 0.10$, and
GraphRAG $0.34 \to 0.16 \to 0.02$, though that last row is radius-limited and is
read with the qualification given below. Think-on-Graph is the exception to the
monotonicity but not to the direction: it falls from $0.53$ to $0.37$ and
partially recovers to $0.45$, still $0.08$ below its own one-hop score. The
distinction is worth keeping. ``Every other system decays'' would be the tidier
sentence, but the table does not say it. A system that gets \emph{better} as the
required chain gets longer is not merely ahead on average. Its profile is the
only one in the table that fits a system which actually traverses the chain.
That rising profile is a stronger form of the multi-hop claim than any aggregate
could carry.

\textbf{Think-on-Graph beats AGR on CWQ's one-hop stratum} ($0.53$ against
$0.46$, $n = 137$). This is reported rather than buried because it makes the
rest of the story more credible, not less. \Cref{sec:ablations} identifies the
mechanism: decomposition costs accuracy on questions that do not need
decomposing. CWQ's \textsf{h1} stratum is full of them.

\textbf{Vector-RAG's WebQSP profile is the paradigm in two cells.} It scores
$0.82$ at one hop --- the best single-hop number anywhere in the table, better
than AGR's $0.75$ --- and $0.12$ at two. When one triple suffices, dense
retrieval over verbalised triples is excellent. When a chain is required, it has
nothing to offer. This is the cell the multi-hop argument rests on. A one-triple
context is chain-incapable by construction rather than by configuration.

\textbf{GraphRAG's WebQSP inversion has two causes. Only one of them was
originally named.} It scores $0.30$ at \textsf{h1} against $0.44$ at
\textsf{h2}. Part of that is hub noise: popular one-hop entities have large
neighbourhoods. The answer is drowned in the reranked window --- or, more
exactly, drowned \emph{or absent}, since the retrieval step samples at most
$100$ edges per topic entity without an ordering, and $55.1\%$ of its topic
entities carry more edges than that (\Cref{sec:baseline-graphrag}). Which of the
two applies to any given question is not recoverable from the run records. But
the inversion also marks the edge of the retrieval radius. The same section
scopes that as a bound on this implementation rather than on static retrieval as
such. WebQSP's \textsf{h2} stratum is largely mediator paths, which a one-hop
expansion reaches by hopping through the mediator. CWQ's \textsf{h2}, by
contrast, is genuine composition, which it cannot reach at all: $0.44$ against
$0.16$ on nominally the same stratum depth. Read GraphRAG's per-stratum row as
measuring what its radius covers, not how far static retrieval can see.

\textbf{The unreachable stratum is a hallucination probe, not an accuracy
measurement.} For these questions, by construction, no path to gold exists
within the hop bound. So every ``hit'' is an assertion the environment cannot
support, one that just happened to match the label. Vector-RAG scores $0.67$ on
CWQ's three unreachable questions and no-retrieval $0.17$ on WebQSP's twelve.
The counts are tiny, and no weight is placed on them. But the direction is
exactly what \Cref{sec:groundedness-results} measures at scale: Hits@1 will
happily reward what RQ2 is about the absence of.

\section{Groundedness and Hallucination Rate Across
Systems}\label{sec:groundedness-results}

\begin{table}[!tb]
  \begin{center}
    \caption[Two-tier groundedness]{Two-tier groundedness. Tier~1 is
    deterministic over every asserted entity in all $4{,}000$ records. Tier~2 is
    a language-model entailment judgement over a stratified sample of $60$
    assertions per system per dataset. The judge itself was validated
    separately, at $\kappa = 0.6995$ over $100$ paired annotations
    (\Cref{sec:judge}).}
    \label{tab:groundedness}
    \small
    \begin{tabular}{|l|l|r|r|r|r|}
      \hline
      & & \multicolumn{3}{c|}{\textbf{Tier 1: structural}} & \textbf{Tier 2} \\
      \hline
      \textbf{Data} & \textbf{System} & \textbf{Entities} & \textbf{Ungrounded}
      & \textbf{Questions} & \textbf{Supported} \\
      \hline
      WebQSP & No-retrieval   & 661  & 179 (27.1\%) & 37.6\% & 63.3\% \\
      WebQSP & Vector-RAG     & 1040 & 38 (3.7\%)   & 6.2\%  & 61.7\% \\
      WebQSP & GraphRAG       & 800  & 6 (0.8\%)    & 2.8\%  & 55.0\% \\
      WebQSP & Think-on-Graph & 832  & \textbf{0 (0.0\%)} & \textbf{0.0\%} &
        61.7\% \\
      WebQSP & AGR            & 1235 & \textbf{0 (0.0\%)} & \textbf{0.0\%} &
        \textbf{66.7\%} \\
      \hline
      CWQ & No-retrieval   & 340 & 42 (12.4\%) & 13.7\% & 36.7\% \\
      CWQ & Vector-RAG     & 289 & 7 (2.4\%)   & 3.4\%  & \textbf{50.0\%} \\
      CWQ & GraphRAG       & 224 & 2 (0.9\%)   & 1.3\%  & 36.7\% \\
      CWQ & Think-on-Graph & 433 & \textbf{0 (0.0\%)} & \textbf{0.0\%} & 43.3\%
        \\
      CWQ & AGR            & 474 & \textbf{0 (0.0\%)} & \textbf{0.0\%} & 48.3\%
        \\
      \hline
    \end{tabular}
  \end{center}
\end{table}

\subsection{Tier 1: Structural Hallucination Is a Property of the Paradigm}

The differential is emphatic. The parametric control asserts $661$ entities
on\index{hallucination}\index{groundedness} WebQSP, of which $179$ --- $27.1\%$
--- have no basis in the knowledge environment. And $37.6\%$ of the questions it
answers contain at least one such entity. Both navigating systems sit at exactly
zero, over $1{,}235$ and $832$ asserted entities. Parametric answering
hallucinates structurally at double-digit rates. Graph-navigated answering
removes structural hallucination outright.

\textbf{But Think-on-Graph reaches zero as well. That changes the attribution.}
Structural groundedness is a property of the \emph{navigation paradigm}. Any
system that copies its answer entities out of triples it actually traversed
inherits it. It is not a property of this thesis's verification layer
specifically. So the honest claim is narrower, and more testable, than the one
the work set out to make: \emph{both navigating systems achieve zero structural
hallucination; they differ at the level of semantic support}. Think-on-Graph's
precision of $0.571$ and $0.403$, against AGR's $0.697$ and $0.484$
(\Cref{tab:main}), shows it asserting far more entities that exist, connect to
the topic, were copied from real triples --- and are not the answer.
\Cref{sec:limitations} builds on exactly this gap.

Two readings of the middle band belong on the record. The no-retrieval control's
CWQ rate ($12.4\%$) is less than half its WebQSP rate ($27.1\%$). That is not
because it fabricates less on harder questions. It is because it hedges more:
$38.0\%$ against $12.2\%$. Abstention pushes the measured hallucination rate
down. That is the abstention trade appearing in a third metric. And Vector-RAG's
$38$ ungrounded entities are mostly not fabrications: its index is global. So
retrieval can surface facts about a similarly-named subject that exists in the
graph but has no connection to the question's topics. The model then copies
faithfully from them. The metric correctly flags those as unanchored to the
question. But the mechanism is retrieval error rather than generation error.
Structural ungroundedness conflates the two. The distinction is recoverable from
whether the entity appears in the system's own retrieved facts.

\subsection{Tier 2: A Narrow Band and an Underpowered Comparison}

Tier~2 asks whether the asserted entity is supported \emph{as the
answer}.\index{Cohen's kappa} Under the instruction that mere connectedness is
not support, every system lands in a $36.7$--$66.7\%$ band. AGR is highest on
WebQSP at $66.7\%$. On CWQ it is second: $5.0$ points ahead of Think-on-Graph,
which is the comparison the rest of this chapter makes, but $1.7$ behind
Vector-RAG's $50.0\%$. Its margin over the agentic baseline is $5.0$ points on
both datasets. Its lead over \emph{every} system is a WebQSP result only.

Neither gap is resolvable at this sample size. With $n = 60$ per cell the $95\%$
binomial interval has a half-width of roughly $12$--$13$ points. So a $5$-point
difference is inside the noise --- as is Vector-RAG's $1.7$. No conclusion in
this thesis rests on either. What the tier does establish is the level: semantic
support is \emph{hard} for every system, including the two that are structurally
perfect. Zero structural hallucination coexists with roughly a third to a half
of assertions failing a strict semantic-support test. That coexistence is the
single most useful thing the metric says. It says it about the paradigm, not
about any one system.

Two properties of the measurement limit how it can be read, and both were
pre-specified. The rates are conservative lower bounds, depressed by the same
amount across systems. That uniformity keeps the between-system comparison
intact and makes the absolute levels untrustworthy. And the judge was validated
at $\kappa = 0.6995$ against a pre-specified bar of $0.7$ (\Cref{sec:judge}):
substantial agreement, just short of the threshold. Both are reasons to treat
this tier as corroborating rather than decisive.

\section{Efficiency and Cost}\label{sec:efficiency}

\subsection{Token and Call Budgets per System}

\begin{table}[!tb]
  \begin{center}
    \caption[Cost per question]{Cost per question. The seconds column uses
    cold-cache records only. \textbf{Clipped} is the share of questions on which
    the shared $25$-call budget ran out mid-search.}
    \label{tab:efficiency}
    \small
    \begin{tabular}{|l|l|r|r|r|r|}
      \hline
      \textbf{Data} & \textbf{System} & \textbf{Tokens} & \textbf{Calls} &
      \textbf{Seconds} & \textbf{Clipped} \\
      \hline
      WebQSP & No-retrieval   & 113   & 1.0  & 1.2  & --- \\
      WebQSP & Vector-RAG     & 527   & 1.0  & 1.4  & --- \\
      WebQSP & GraphRAG       & 531   & 1.0  & 2.1  & --- \\
      WebQSP & Think-on-Graph & 3{,}615 & 12.8 & 12.9 & 29.2\% \\
      WebQSP & AGR            & 4{,}511 & 6.2  & 11.2 & 0.0\% \\
      \hline
      CWQ & No-retrieval   & 118   & 1.0  & 1.0  & --- \\
      CWQ & Vector-RAG     & 535   & 1.0  & 1.4  & --- \\
      CWQ & GraphRAG       & 470   & 1.0  & 2.2  & --- \\
      CWQ & Think-on-Graph & 5{,}572 & 18.2 & 15.2 & 44.0\% \\
      CWQ & AGR            & 6{,}818 & 8.9  & 17.2 & 0.0\% \\
      \hline
    \end{tabular}
  \end{center}
\end{table}

Agency costs roughly ten times what one-shot retrieval costs: $4{,}511$
and\index{budget} $6{,}818$ tokens per question against $\sim 500$, and
$11$--$17$ seconds against one or two. That is the price of RQ1's accuracy, and
it is stated as a price.

The informative comparison is with the other agentic system, because the two run
under the same $25$-call cap. AGR uses \emph{half} the calls --- $6.2$ against
$12.8$ on WebQSP, $8.9$ against $18.2$ on CWQ --- with more tokens per call,
because each call carries a richer prompt: a plan, a scored candidate set, a
fact window. Fewer, better-informed decisions rather than more, blinder ones.

The clip column is the sharpest form of that. Beam search costs are
multiplicative in width and depth. So under the shared budget Think-on-Graph
exhausts its calls before finishing on $29.2\%$ of WebQSP questions and $44.0\%$
of CWQ questions, while AGR never does. On nearly half the multi-hop benchmark
the baseline is answering from a truncated search. This is a fair constraint
(\Cref{sec:baseline-tog}): both systems live under the same cap. It is the
efficiency finding rather than a caveat on it: under an identical budget, AGR
converts its calls into correct answers at roughly twice the rate.

Never reaching the call cap is not the same as running unbounded. The
distinction matters. The rest of this thesis restates this result. AGR's own
budgets bind, the depth cap of \Cref{tab:budget} hardest: the meter refuses a
deeper expansion on $25.6\%$ of test questions, against $6.2\%$ for backtracks
and $1.9\%$ for repair iterations. \Cref{sec:repair} and
\Cref{sec:verify-example} each narrate a run that ends that way. What the clip
column establishes is the narrower and more useful thing --- that of the two
systems run under one shared call budget, that budget truncates only one of
them.

\subsection{The Accuracy--Cost Frontier}

\begin{figure}[!tb]
  \begin{center}
    % GENERATED by scripts/build_figures.py from
% results/phase4/thesis_numbers.json -- do not edit.

\begin{tikzpicture}
\begin{axis}[agrplot, name=plot1, 
  title={WebQSP},
  xmode=log, log basis x=10,
  xlabel={Mean tokens per question}, ylabel={Hits@1},
  xmin=80, xmax=12000, ymin=0.15, ymax=0.85,
  legend style={at={(1.08,-0.30)}, anchor=north, legend columns=5, column sep=1.8mm},
  legend cell align=left,
  legend image post style={solid},
]
\addplot[only marks, mark=o, mark size=2.6pt, color=black!55] coordinates {(113,0.453)};
\addlegendentry{No-retrieval}
\addplot[only marks, mark=square, mark size=2.6pt, color=black!70] coordinates {(527,0.568)};
\addlegendentry{Vector-RAG}
\addplot[only marks, mark=triangle, mark size=2.6pt, color=black!70] coordinates {(531,0.34)};
\addlegendentry{GraphRAG}
\addplot[only marks, mark=diamond, mark size=2.6pt, color=black!85] coordinates {(3615,0.66)};
\addlegendentry{Think-on-Graph}
\addplot[only marks, mark=*, mark size=2.6pt, color=black] coordinates {(4511,0.755)};
\addlegendentry{AGR}
\end{axis}
\begin{axis}[agrplot, name=plot2, at={($(plot1.east)+(12mm,0)$)}, anchor=west,
  title={ComplexWebQuestions},
  xmode=log, log basis x=10,
  xlabel={Mean tokens per question}, ylabel={Hits@1},
  xmin=80, xmax=12000, ymin=0.15, ymax=0.85,
%
]
\addplot[only marks, mark=o, mark size=2.6pt, color=black!55] coordinates {(118,0.307)};
\addplot[only marks, mark=square, mark size=2.6pt, color=black!70] coordinates {(535,0.203)};
\addplot[only marks, mark=triangle, mark size=2.6pt, color=black!70] coordinates {(470,0.205)};
\addplot[only marks, mark=diamond, mark size=2.6pt, color=black!85] coordinates {(5572,0.435)};
\addplot[only marks, mark=*, mark size=2.6pt, color=black] coordinates {(6818,0.522)};
\end{axis}
\end{tikzpicture}

  \end{center}
  \caption[Accuracy against token cost]{Accuracy against token cost, one point
  per system. The horizontal axis is logarithmic: the agentic systems cost
  roughly an order of magnitude more than the one-shot ones, and a linear axis
  would compress every non-agentic system into a single column. Read this figure
  for the token frontier only. Think-on-Graph sits below \emph{and} to the left
  of AGR, so it is a frontier point on this axis, not a dominated one. The
  domination claim holds on calls, which are not plotted and are given in
  \Cref{tab:efficiency}. Generated from
  \textsf{results/phase4/thesis\_numbers.json}.}
  \label{fig:accuracy-cost}
\end{figure}

Reading \Cref{tab:main} against \Cref{tab:efficiency}, and plotted
as\index{budget} \Cref{fig:accuracy-cost}, the frontier depends on which cost is
charged. The two costs this thesis records do not agree.

\textbf{On tokens, the axis of \Cref{fig:accuracy-cost}, Think-on-Graph is not
dominated.} It buys its lower accuracy at a genuinely lower token price on both
datasets --- $3{,}615$ against AGR's $4{,}511$ on WebQSP, $5{,}572$ against
$6{,}818$ on CWQ --- and sits on the Pareto set of both, alongside the
no-retrieval control at $113$ tokens and, on WebQSP, Vector-RAG at $527$. The
interior points are the static graph retrievers, which the parametric control
beats on both axes at once. Any claim that Think-on-Graph is dominated has to
name an axis on which that is true.

\textbf{On calls it is dominated. Calls are what the budget meters}
(\Cref{sec:budgets}). Think-on-Graph spends $12.8$ and $18.2$ calls against
AGR's $6.2$ and $8.9$ --- lower Hits@1 at roughly double the call cost, on both
datasets, with no interior ambiguity. The two costs diverge because AGR spends
more than twice as many tokens per call, converting fewer, larger calls into a
comparable token total. Which cost binds is a deployment question: token billing
separates the two systems very little, whereas a per-call rate or latency budget
separates them by a factor of two. It is the call budget that clips
Think-on-Graph on $29\%$ and $44\%$ of questions (\Cref{sec:tog-split}).

Around that, the corners are unchanged. Vector-RAG on WebQSP is the cheap
corner, $0.568$ Hits@1 for $527$ tokens and $0.82$ on the one-hop stratum, which
makes it the right choice for shallow questions and useless for deep ones. AGR
is the accurate corner on both datasets. And the no-retrieval control holds the
cheap end of the token frontier while being off any accuracy-worthy frontier
once precision is counted: it remains the cheapest way to be right about half
the time on WebQSP, which is the contamination point restated as an economic
one.

\section{Ablation Study}\label{sec:ablations}

Four single-deviation conditions on stratified halves ($n = 200$ WebQSP, $n =
198$ CWQ), each against the unmodified system restricted to the same questions.
\Cref{tab:ablation} reports the conditions. \Cref{tab:rq3} summarises the
component attribution. Throughout, $\Delta$ is \emph{ablated minus reference}.
So a positive $\Delta$ means removing the component \emph{improved} the metric.

\begin{table}[!tb]
  \begin{center}
    \caption[Ablation conditions on stratified halves]{Ablation conditions on
    stratified halves. The reference row is the frozen system restricted to the
    same questions, at no additional cost. P and R are the per-question means of
    \Cref{sec:metrics}, reported for the full systems in \Cref{tab:main} and
    carried here so that each condition's precision effect can be read off
    directly. Cost is given in tokens and calls only. There is no seconds
    column, because every record in the \textsf{no-verifier} condition is a full
    cache replay and latency would not be comparable across conditions.}
    \label{tab:ablation}
    % Nine columns: the seven of \Cref{tab:main} plus the two cost columns. Two
    % extra columns of bracketed intervals will not seat at \small however far
    % the padding is trimmed, so this table alone drops a size.
    \footnotesize
    \setlength{\tabcolsep}{3pt}
    \begin{tabular}{|l|l|l|l|r|r|r|r|r|}
      \hline
      \textbf{Data} & \textbf{Condition} & \textbf{Hits@1} & \textbf{F1} &
      \textbf{P} & \textbf{R} & \textbf{Hedge} & \textbf{Tok.} & \textbf{Calls}
      \\
      \hline
      WebQSP & reference       & 0.755 [0.695, 0.810] & 0.659 [0.599, 0.717] &
        0.710 & 0.664 & 8.5\%  & 4{,}562 & 6.2 \\
      WebQSP & no-planner      & \textbf{0.830} [0.775, 0.880] & \textbf{0.742}
        [0.687, 0.794] & 0.788 & 0.749 & 5.5\%  & 3{,}159 & 4.0 \\
      WebQSP & no-backtrack    & 0.745 [0.685, 0.805] & 0.646 [0.586, 0.705] &
        0.704 & 0.649 & 9.5\%  & 4{,}261 & 5.8 \\
      WebQSP & no-verifier     & 0.760 [0.700, 0.820] & 0.663 [0.604, 0.721] &
        0.711 & 0.673 & 8.0\%  & 4{,}460 & 6.1 \\
      WebQSP & embedding-only  & 0.740 [0.680, 0.800] & 0.643 [0.583, 0.702] &
        0.701 & 0.644 & 13.5\% & 3{,}413 & 4.4 \\
      \hline
      CWQ & reference       & 0.495 [0.424, 0.566] & 0.445 [0.379, 0.511] &
        0.461 & 0.451 & 23.2\% & 7{,}105 & 9.1 \\
      CWQ & no-planner      & 0.434 [0.369, 0.505] & 0.398 [0.333, 0.464] &
        0.407 & 0.405 & 26.8\% & 5{,}617 & 6.7 \\
      CWQ & no-backtrack    & 0.490 [0.419, 0.561] & 0.445 [0.379, 0.511] &
        0.456 & 0.453 & 28.3\% & 5{,}792 & 7.4 \\
      CWQ & no-verifier     & 0.490 [0.419, 0.561] & 0.436 [0.371, 0.503] &
        0.451 & 0.446 & 20.2\% & 6{,}704 & 8.7 \\
      CWQ & embedding-only  & 0.475 [0.409, 0.545] & 0.437 [0.370, 0.502] &
        0.452 & 0.444 & 28.3\% & 4{,}925 & 5.9 \\
      \hline
    \end{tabular}
  \end{center}
\end{table}

\subsection{Without the Planner: A Stratum-Dependent Effect}

This is the one condition that reaches significance. And it points the
opposite\index{planner} way from what the design intended.

On WebQSP, \emph{removing} the planner improves every axis at once: Hits@1
$0.755 \to 0.830$, F1 $0.659 \to 0.742$, hedge rate $8.5\% \to 5.5\%$, tokens
$4{,}562 \to 3{,}159$, calls $6.2 \to 4.0$. McNemar gives $p = 5.9 \times
10^{-3}$ on $6$ against $21$ discordant pairs. The per-stratum breakdown locates
the damage precisely: \textsf{h1} rises $0.75 \to 0.85$ and \textsf{h2} rises
$0.84 \to 0.89$. So on this dataset decomposition hurts simple \emph{and}
two-hop questions. The planner prompt's first worked example tells the model not
to decompose single-hop questions. The ablation proves that instruction fires
too rarely, and it proves it with a $p$-value.

On CWQ the sign reverses. Significance is not reached: Hits@1 $0.495 \to 0.434$,
$27$ against $15$ discordant pairs, $p = 0.088$. The stratum table shows a
coherent mechanism rather than noise. \textsf{h1} \emph{improves} ($0.49 \to
0.54$), \textsf{h2} collapses ($0.48 \to 0.34$), and \textsf{h3plus} moves
little ($0.62 \to 0.54$). That is exactly the pattern ``decomposition helps
genuine composition and hurts questions that do not need it''. It does not clear
$\alpha = 0.05$ at $n = 198$. The honest statement carries both halves:
\emph{the planner's benefit depends on the stratum and is significant on WebQSP
($p = 0.006$); the CWQ pattern points the same way but does not reach
significance at this sample size ($p = 0.088$)}.

So the finding is neither ``the planner helps'' nor ``the planner hurts''. The
planner's value depends on question structure. Applying decomposition to every
question regardless is the defect. \Cref{sec:decomposition-errors} reads the
$21$ and $15$ questions where removing the planner flipped a failure into a
success and names the mechanisms. \Cref{sec:future} turns them into an adaptive
gating proposal that this ablation earned rather than motivated.

\subsection{Without Backtracking}

Hits@1 moves $-0.010$ on WebQSP and $-0.005$ on CWQ. F1 moves $-0.013$ and
$0.000$.\index{backtracking} Discordant pairs are $5$ against $3$ and $6$
against $5$, giving $p = 0.73$ and $p = 1.00$. There is no detectable effect at
this sample size.

Two readings are available. The data does not separate them: backtracking's
contribution may be small because the scorer and the verifier already prevent
most of the excursions it exists to correct, or $n \approx 200$ may simply be
too few questions to resolve a small effect. This is reported as ``no detectable
effect at $n \approx 200$'', never as a confirmed null. The mechanism is not
idle --- the full runs record $63$ backtracks on WebQSP and $199$ on CWQ,
dominated by evaluator-triggered ones ($47$ and $170$). So the condition removes
real activity whose accuracy consequence this experiment cannot measure.

\subsection{Without the Verification Layer}

The cleanest null in the table, and the most informative one. Discordant pairs
are\index{verification layer} $0$ against $1$ on WebQSP and $1$ against $0$ on
CWQ: across $398$ paired questions, removing claim verification changed the
correctness of exactly two. Hits@1 moves $+0.005$ and $-0.005$. F1 moves
$+0.004$ and $-0.009$. Both $p = 1.00$.

This reproduces the development-set result of \Cref{sec:dev-tuning} exactly.
\Cref{sec:verification-layer} predicted as much, by construction. On $77$ of
$80$ development questions the layer certified the draft and left it alone; the
repair path never ran. A mechanism that mostly passes cannot move an accuracy
metric.

The conclusion this licenses is narrower and sharper than ``the verifier
helps''. It is the one \Cref{sec:findings} adopts: \emph{claim verification does
not measurably change how often AGR is right; it changes how often AGR is right
for a reason it can exhibit}. Hits@1 cannot see that distinction. What makes it
checkable is the supporting-triple output contract of
\Cref{sec:output-contract}. The evidence for the layer is that contract and the
specimen analyses of \Cref{sec:verifier-errors}. It is neither of the two
aggregate results it would be natural to reach for instead.
\Cref{tab:groundedness} is reached by the agentic baseline as well. That is why
\Cref{sec:findings} attributes it to graph navigation. And the precision column
of \Cref{tab:main} does not move when the layer is removed, as the P column of
\Cref{tab:ablation} shows directly. It is not this row either. This row is
reported at full length rather than omitted.

\subsection{Embedding-Only Scoring}\label{sec:abl-embonly}

Removing the language-model term from the hybrid scorer costs $-0.015$
and\index{hybrid scoring function} $-0.020$ Hits@1 and $-0.016$ and $-0.008$ F1
--- consistently negative on both datasets and both metrics, and consistently
short of significance ($p = 0.66$ and $p = 0.48$).

Two corroborating details make the direction credible without over-claiming from
it. The condition is cheaper than the reference on both datasets --- $3{,}413$
against $4{,}562$ and $4{,}925$ against $7{,}105$ tokens, $4.4$ against $6.2$
and $5.9$ against $9.1$ calls. So the language-model scoring term is buying
accuracy with tokens as designed. And it carries by far the most backtracks,
$57$ and $143$ against the reference's $38$ and $108$, with the
\textsf{low\_score} trigger that $\tau$ governs roughly tripling in both
datasets and accounting for about three quarters of the WebQSP increase and
about half of CWQ's.

That second detail was originally read as the $\tau$ mechanism of
\Cref{sec:backtracking} visible in the logs. The reading survives only partly.
The signal $\tau$ thresholds is not computed identically in the two conditions,
for the reason \Cref{sec:backtracking} records: at $\alpha = 1$ the two
auxiliary maxima are never taken. So $\sigma$ is a maximum over beam-selected,
non-banned edges rather than over every candidate offered. After a backtrack has
banned the best relation, $\sigma$ falls in this condition, where it holds
steady in the reference. That divergence adds \textsf{low\_score} backtracks
through an implementation asymmetry, not through the scoring change under test.
The two effects push the same way. Separating them would mean re-running the
condition with the auxiliary terms restored, and that would no longer be the
frozen system.

So the honest statement is narrower than the original one. Removing the model
term raises the backtrack count, but an unknown and non-trivial share of that
rise is an artefact of how $\sigma$ is computed when the term is disabled. The
backtrack count is corroborating, not load-bearing.

Note which way the defect cuts. Extra backtracks make the embedding-only
condition look busier than the scoring change alone would, so the defect
flatters the component under test rather than the null. Nothing in
\Cref{tab:rq3} depends on it, since the accuracy and F1 deltas come from
answers, not from backtracks. The main evidence for $\alpha = 0.7$ remains the
development sweep of \Cref{tab:sweep}, where the $\alpha = 1.0$ condition lost
five hits on $80$ questions. This ablation corroborates it at test scale rather
than independently establishing it.

\subsection{Paired Significance Testing and Statistical Power}

\begin{table}[!tb]
  \begin{center}
    \caption[Component attribution]{Component attribution, RQ3. $\Delta$F1 is
    \emph{ablated minus reference}, so positive means removing the component
    improved F1. Only the planner condition reaches significance.}
    \label{tab:rq3}
    \small
    \begin{tabular}{|l|l|l|l|}
      \hline
      \textbf{Component} & \textbf{WebQSP $\Delta$F1 ($p$)} & \textbf{CWQ
      $\Delta$F1 ($p$)} & \textbf{Verdict} \\
      \hline
      Planner & $+0.083$ ($0.006$) & $-0.047$ ($0.088$) & Stratum-dependent \\
      Backtracking & $-0.013$ ($0.73$) & $0.000$ ($1.00$) & No detectable effect
        \\
      Verification & $+0.004$ ($1.00$) & $-0.009$ ($1.00$) & No accuracy effect
        \\
      $\alpha$-blend & $-0.016$ ($0.66$) & $-0.008$ ($0.48$) & Directional,
        underpowered \\
      \hline
    \end{tabular}
  \end{center}
\end{table}

One of four conditions reaches significance. That is the correct outcome to
report, not a disappointing one. A correction for multiple comparisons was not
among the pre-specified decisions of \Cref{sec:pre-registration}. Applying one
now would be a choice made after seeing the results. So the uncorrected
$p$-values are reported as computed and read as descriptive rather than as a
family-wise claim. \Cref{app:power} gives the arithmetic behind the second
qualification: at $n \approx 200$ McNemar's test needs a discordance pattern
this experiment could not have produced for the effects it was looking for. So
three conditions are reported as ``no detectable effect at this sample size''
rather than as no effect. The distinction is the whole content of
\Cref{sec:threats}'s entry on underpowered ablations.

\section{Error Analysis}\label{sec:erroranalysis}

\Cref{sec:results} measured how often each system fails. This analysis reads the
failures and names the mechanism of each one. It is the part of the work no
script produced: $256$ failure packets --- question, gold, answer, plan,
explored relations with their scores, backtrack reasons, verifier outcome ---
read and labelled by hand, one category and one sentence each.

The reading protocol was fixed before it began. For each packet, find where the
trajectory \emph{first} left the correct path, and label that point rather than
the last visible symptom. Where two causes compete, take the earlier plausible
one and record the other. Without that rule, a run that mis-decomposes, then
explores the wrong relation, then hedges would be counted three times under
three categories.

A note on the identifiers used throughout this analysis, because they are meant
to be looked up. A ComplexWebQuestions identifier is the WebQSP identifier it
was derived from followed by a thirty-two character hash. The stem alone does
not identify a question: within these samples one stem covers up to twelve
distinct CWQ questions. Thirteen stems are themselves complete WebQSP
identifiers naming something else entirely. Cases are therefore cited as stem
plus the first eight hash characters. That combination is unique within both
samples. \Cref{app:samples} lists every identifier in full.

\subsection{Annotation Protocol and the Failure Taxonomy}\label{sec:taxonomy}

Ten categories carry the census, each naming a distinct architectural locus
rather than a symptom, so that a count against a category points at the
component that produced it. \Cref{app:taxonomy-schema} gives the schema in full,
together with the open subtype vocabulary beneath it --- which drives nothing
and exists so that a mechanism can be named before it is understood --- and the
one naming decision that was got wrong and corrected.

\paragraph{Population, Not Sample}

The census\index{failure census} is a \emph{population}, not a sample: every
question the full system did not answer correctly was read and labelled. Three
sets are set aside before it is taken, of which only the adjudicated benchmark
defects of \Cref{sec:gold-quality} remove questions from the analysis
altogether. The other two are of questions that were never failures.
\Cref{app:census-denominator} reconciles the three, which overlap, and shows how
the $256$ read failures follow. That the census is complete rather than sampled
is what licenses the category counts below being read as counts rather than as
estimates.

\subsection{Distribution of Failure Categories Across
Datasets}\label{sec:failure-distribution}

% This table and the figure below hold the same counts in two notations and are
% meant to be read together, so they are allowed onto a float page (p) rather
% than being split across two text pages.
\begin{table}[!tbp]
  \begin{center}
    \caption[The merged failure histogram]{The merged failure histogram, $n =
    256$. Wrong and hedge are kept separate. Counts come from
    \textsf{scripts/synthesize\_census.py} over the three label sources of
    \Cref*{tab:census-pop}. The table body is generated from
    \textsf{results/phase4/thesis\_numbers.json}.}
    \label{tab:histogram}
    \small
    % !TEX root = ../buetcsepgthesis.tex
% GENERATED by scripts/build_figures.py from
% results/phase4/thesis_numbers.json -- do not edit.

\begin{tabular}{|l|r|r|r|r|r|}
\hline
& \multicolumn{2}{c|}{\textbf{WebQSP}} &
  \multicolumn{2}{c|}{\textbf{CWQ}} & \\
\hline
\textbf{Category} & \textbf{wrong} & \textbf{hedge} & \textbf{wrong} &
\textbf{hedge} & \textbf{Total} \\
\hline
relation\_selection    &  20 &   6 &  20 &  19 & \textbf{65} \\
composite\_claim       &   1 & --- &  23 &  23 & \textbf{47} \\
kg\_gap                &   8 &   4 &   8 &  24 & \textbf{44} \\
decomposition\_error   &  15 &  11 &  10 &   2 & \textbf{38} \\
answer\_selection      &   1 &   1 &   8 &   5 & \textbf{15} \\
echo                   &   5 &   2 &   4 &   2 & \textbf{13} \\
ambiguous\_question    &   1 &   1 &   3 &   5 & \textbf{10} \\
premature\_termination &   1 &   4 &   1 &   2 & \textbf{8} \\
gold\_noise            &   1 & --- &   3 &   3 & \textbf{7} \\
other                  & --- &   1 & --- &   5 & \textbf{6} \\
verifier\_fn           & --- &   3 &   1 &   1 & \textbf{5} \\
verifier\_fp           & --- & --- & --- &   1 & \textbf{1} \\
\hline
\textbf{Total} & \textbf{53} & \textbf{33} & \textbf{81} & \textbf{92} & \textbf{259} \\
\hline
\end{tabular}

  \end{center}
\end{table}

\paragraph{The Shape Flip, and What It Is Not}

The two datasets fail differently, but not in the way a first reading of
\Cref{tab:histogram} suggests. \textsf{relation\_selection} sits at or next to
the top on \emph{both} --- $25$ of WebQSP's $85$, one behind
\textsf{decomposition\_error}, and $39$ of CWQ's $171$, behind
\textsf{composite\_claim} alone. So it is not what distinguishes them. Picking
the wrong edge is the base rate of graph navigation. It is dataset-independent.

What distinguishes them is a pair of categories that are nearly absent from one
and dominant in the other. \textsf{composite\_claim} is $1$ of $85$ on WebQSP
and $46$ of $171$ on CWQ. \textsf{kg\_gap} is $12$ and $31$. In the other
direction, \textsf{decomposition\_error} is $26$ on WebQSP --- the near-tie
noted above, the two together accounting for $60\%$ of its failures --- against
$12$ on CWQ.

That asymmetry is demonstrable rather than assertable, from two facts already
established. First, the stratum distribution of \Cref{tab:testsets}: WebQSP is
$64\%$ one-hop with a four-question deep tail. So most of its questions have no
composition to get wrong and no constraint to drop, while presenting the planner
with exactly the questions that do not need decomposing. That mismatch is also
the significant ablation result of \Cref{sec:ablations}, appearing here as a
category count. Second, the expressiveness limits certified in
\Cref{sec:unanswerable}: the environment carries no date literals, no numeric
literals, and no ordinals. CWQ's SPARQL templates invoke all three constantly,
while WebQSP's naturally-occurring questions rarely do. The failure profile of a
system on a dataset is, to this extent, a property of the dataset's
construction.

\paragraph{Hedging Where It Should}

A second reading of \Cref{tab:histogram} is worth stating as a positive
result\index{hedge} rather than only as a deficit. On CWQ, \textsf{kg\_gap}
skews three-to-one toward hedges --- $24$ against $8$ --- while on WebQSP it
leans the other way, $4$ against $8$. When CWQ's numeric identifiers, date
literals, and ordinal constraints are unreachable, the system much more often
declines than asserts something wrong. That is the behaviour the anti-vacuity
and grounding provisions of \Cref{sec:claim-decomposition} were written to
produce. It is visible in the one place where the environment guarantees the
system cannot succeed.

\section{Failures in Planning, Navigation, and
Verification}\label{sec:failure-categories}

The categories in this section each sit inside one component of the agent: the
planner and the drafter, the relation scorer, and the verifier. The three
subsections read them in that order, and a count against any of them points at
the code that produced it.

\subsection{Decomposition, Drafting, and Answer-Selection
Errors}\label{sec:decomposition-errors}

Three families share this section because they share a locus: everything
downstream of retrieval and upstream of the answer. By the time these failures
happen, the reasoning is often already complete and correct. That timing makes
them diagnostic, and it makes them cheap to fix.

\paragraph{The Entity-Extraction Bug}

Nine of the $38$ \textsf{decomposition\_error} cases carry the subtype
\textsf{extraction\_bug}, and they are the most actionable finding in the
census. On profession-shaped and ``what did $X$ do'' questions, the evaluator
resolves every correct gold value and the drafted text states them word for
word. Then \textsf{answer\_entities} collapses to the grammatical \emph{subject}
of the sentence. \Cref{tab:extraction} gives three instances read from the raw
traces.

\begin{table}[!tb]
  \begin{center}
    \caption[The entity-extraction bug]{The entity-extraction bug. In each case
    the drafted answer text names every gold value correctly, and the scored
    entity list holds only the sentence's grammatical subject.}
    \label{tab:extraction}
    \small
    \begin{tabular}{|l|r|c|l|}
      \hline
      \textbf{Question ID} & \textbf{Gold} & \textbf{Named in text} &
      \textbf{Scored entities} \\
      \hline
      WebQTest-1215 & 6 professions & all 6 & \textsf{[Stephen Covey]} \\
      WebQTest-704  & 6 professions & all 6 & \textsf{[Thor Heyerdahl]} \\
      WebQTrn-124\_0782789f   & 4 films       & all 4 & \textsf{[Angelina
        Jolie]} \\
      \hline
    \end{tabular}
  \end{center}
\end{table}

\Cref{tab:extraction} understates how complete the drafts are. For
\textsf{WebQTest-1215}, ``who was stephen r covey'', gold lists Author, Manager,
Writer, Consultant, Professor, and Motivational speaker. The drafted answer
reads \emph{``Stephen Covey was a Writer, Motivational speaker, Author,
Consultant, Professor, and Manager''}: all six, correct, in a different order.
\textsf{answer\_entities} comes back as \textsf{['Stephen Covey']}. Navigation
found the answers, the verifier certified them, and the prose reports them. The
step that fails is the one that turns answer text into an entity list, and it
reaches for the subject rather than the predicate values every time.

The consequence is a limitation on \Cref{sec:results}, not only a bug. Scoring
reads \textsf{answer\_entities}, so \Cref{tab:main} records every one of these
cases as a wrong answer or a hedge. On ``list the professions'' and ``list what
$X$ did'' question shapes, an extraction defect that has nothing to do with
reasoning quality pulls the reported accuracy down. Nine cases in $256$ cannot
move the headline numbers much. But the measured figure for that question shape
is a floor, and it should be read as one. \Cref{sec:future} lists the fix among
the repairs the census earned.

\paragraph{Right Candidate Retrieved, Wrong One
Drafted}\label{sec:answer-selection}

\textsf{answer\_selection} ($15$ cases, $13$ of them on CWQ) is the extraction
bug's more serious sibling. Here the exact gold value demonstrably entered the
resolved candidate set, and the drafter chose something else.
\textsf{WebQTest-1555} asks what the parliament of Nepal is called. The
evaluator resolves \textsf{Parliament of Nepal} itself; the draft names its two
component chambers instead. \textsf{WebQTrn-2784\_abedc8ac} asks which Tupac
Shakur film Malcolm Campbell edited. The trace records
\textsf{resolved=['Nothing but Trouble']}, the literal gold value, which
satisfies both constraints. The draft substitutes \textsf{Gang Related} and then
hedges on it.

One case in this family was first filed as a verifier error, and the correction
is worth giving in full because it is the kind of mistake this taxonomy exists
to prevent. \textsf{WebQTrn-1294\_a4b2006a} was labelled \textsf{verifier\_fn}
on the theory that the verifier had checked a claim about a different candidate
than the one the evaluator resolved. The run record refutes that theory. The
draft itself reads ``Marlon Brando also played in \emph{Joy}'', and the verifier
rejected exactly what the drafter asserted. Brando was not in \emph{Joy}. De
Niro was, and the evaluator had resolved him. \textbf{The verifier was correct.}
The error sits upstream, in the answerer, and the label now reads
\textsf{answer\_selection}. The case is also the clearest example in the census
of the verification layer doing its job. \Cref{sec:verifier-errors} opens by
pointing back to it.

\paragraph{Failing to Commit: Six Cases of Hedging on a Verified Fact}

The \textsf{other} category was left open for mechanisms the schema did
not\index{hedge} anticipate. All six of its cases turned out to share one.
Nothing predicted it, which is why the schema lacks it: \emph{the evaluator
resolves the exact gold value, the verifier returns} \textsf{grounded},
\emph{and the drafted text hedges anyway}. No rejection is involved, so this is
not a verifier error at all.

\textsf{WebQTest-689} states ``Spanish Language was spoken in Spain'' and then,
in the same answer, that the first language spoken in Spain could not be
determined. \textsf{WebQTrn-125\_003c6a04} resolves \textsf{Memphis}, matching
gold exactly, and the draft says the location could not be determined.
\textsf{WebQTrn-1392\_d372995c} names both of Eleanor Roosevelt's schools
correctly and then declares the school undeterminable.
\textsf{WebQTest-989\_4b6636a0} (Battle of Dunkirk),
\textsf{WebQTest-1923\_a64ef0f5} (\emph{The Last Song}), and
\textsf{WebQTrn-2721\_2d207ed9} (\textsf{Alois Hitler}) have the same shape. In
every case the entity list comes back empty and the question scores as a hedge.

The contrast with the extraction bug is exact and worth holding on to. There,
the text is right and the entity list is wrong. Here, the text contradicts a
fact it has just stated. Six instances out of six share the mechanism, and that
uniformity makes it a defect rather than a distribution.

\paragraph{Context-Stripping: When Decomposition Discards the Question}

Subtyping is thinner in this category than elsewhere, and the shortfall
belongs\index{planner} on the record before the cases are read.
\textsf{decomposition\_error} holds $38$ failures. The four named subtypes,
\textsf{extraction\_bug} ($9$), \textsf{paraphrase\_drift} ($9$),
\textsf{context\_stripping} ($3$), and \textsf{over\_decomposition} ($1$), cover
$22$ of them. The other $16$ carry a category label only. So the mechanisms
below are the mechanisms of the $22$, not of all $38$. The subtype vocabulary
stayed open rather than being stretched to cover the residue.

\textsf{paraphrase\_drift} and \textsf{over\_decomposition} are the expected
decomposition failures --- a rewritten sub-objective scoring worse against the
target relation than the raw question text, or a redundant step resolving an
already-unambiguous entity. \textsf{WebQTest-1829} is the clean
over-decomposition case: the plan splits ``what type of religion did
massachusetts have'' into \textsf{['find Massachusetts', 'find the religion type
associated with \#1']}, burns $14$ of $16$ calls on three evaluator backtracks,
and hedges, while the undecomposed run finds the same fact in three calls with
no backtracks.

\textsf{context\_stripping} ($3$ cases) is the mechanism this analysis would
feature if it could feature only one. It is a failure of decomposition that
decomposition's own design cannot see: splitting a question into steps
\emph{removes} a qualifier that appears only in a later clause. So an early
sub-objective resolves the wrong sense of an ambiguous term before the
disambiguating information ever arrives. In \textsf{WebQTest-1367}, which asks
where Glastonbury, England is, the plan's first step anchors on \textsf{'find
Glastonbury'} alone. The explorer resolves toward the identically-named
Connecticut town. In \textsf{WebQTrn-2615\_48a6ac56}, which asks what
instruments the author of \emph{Fela!} plays, the first sub-objective is scored
with no mention of instruments and resolves the wrong sense of an ambiguous
authorship relation, where the undecomposed run answers all five instruments
correctly in fewer calls.

\textsf{WebQTrn-2570\_d63877a4} is the strongest specimen in the census. Its
trace shows every link of the mechanism. The question --- ``This 33rd president
was the nation's leader during WW2'' --- is a single entity under two
descriptions. The planner emits \textsf{['find the 33rd president', "find the
nation's leader during WW2"]}, with no \textsf{\#1} reference between them,
modelling two descriptors of one entity as though they were a chain. Because
sub-objective one is the active scoring context, the anchor \textsf{World War
II} scores $0.128$--$0.133$ --- the noise floor --- while presidency-hub
relations score $0.46$--$0.54$: decomposition actively \emph{suppressed} the one
anchor that could have discriminated. Objective one then never completes
(\textsf{objective\_done: false} at every round). So objective two never
activates. The WW2 constraint is never used at any point in the search. The
agent spends its depth budget wandering around the presidency hub. It then
drafts a plausible-sounding president: Woodrow Wilson, the 28th, and nowhere
near the war. The undecomposed run answers Harry S.~Truman correctly in three
calls. ``WW2'' stayed in the scoring context. Two further observations belong to
that trace. No ordinal relation appears anywhere in it. So the ``33rd''
constraint played no part either --- consistent with the environment's lack of
ordinal encoding (\Cref{sec:kg-gaps}). And the verifier returned
\textsf{grounded} on the Wilson claim, which \Cref{sec:verifier-errors} takes up
as a boundary case rather than a defect.

Three instances of one mechanism, two of them surfaced by the ablation contrast
and one by ordinary reading, is what makes context-stripping a named finding
rather than an anecdote. It is the concrete evidence behind the adaptive-gating
proposal of \Cref{sec:future}. A fourth case belongs beside them and is filed as
plain \textsf{decomposition\_error}, since no subtype fits it:
\textsf{WebQTest-869}, ``where is ancient phoenician'', whose planner rewrote
the sole sub-objective as \textsf{['find ancient Phoenician']}. So the explorer
re-identified the topic instead of locating it. The run answered ``Phoenicia''
--- the question's own subject, grounded, and useless, against an undecomposed
run that answered ``Ancient Phoenicia was in Lebanon'' in three calls. What
separates it from the three above is \emph{what} the rewrite discarded: they
drop a constraint arriving in a later clause, this one drops the interrogative,
the word that says what is being asked for. The result lands in the graph as an
\textsf{echo} (\Cref{sec:echo}) --- with the ask removed, the nearest
well-grounded entity to a topic is the topic. Both losses are visible to a
planner that compares the rewrite against the original question, which is why
the gating proposal checks the rewrite for the interrogative as well as for
redundancy.

\subsection{Relation-Selection and Navigation Errors}\label{sec:relation-errors}

\textsf{relation\_selection} is the largest category at $64$ and carries no
subtypes, because the mechanism is uniform: a wrong relation was explored, or
the right one was explored and abandoned. Its interest is not in the individual
cases but in the fact that several of them are the \emph{same} gap seen from
different angles --- one relation confusion reproduced across three unrelated
questions, and, most sharply, a CWQ triplet reaching the same entity by three
different routes, none of which ever tries the single relation that answers all
three. Three independent entry points into the same resolver hitting the
identical dead end is the cleanest evidence in the corpus that a relation gap is
\emph{systemic} rather than an artefact of one question's phrasing: it is a
property of the scorer's ranking over that entity's relation set, not of the
question.

\textsf{premature\_termination} is small ($8$ cases) and splits into two
mechanisms that must not be merged.\index{evaluator} Five carry the
\textsf{evaluator} subtype and are the \emph{same} systemic pattern --- the
correct relation is explored with a solid score, as high as $0.731$, and the
evaluator backtracks away from it and never returns, across entirely unrelated
domains. A relation scoring $0.73$ being explored and discarded is not a
reasoning gap. It reads as a threshold or backtracking-policy defect. It is one
of the most directly fixable findings in the census. The other three carry
\textsf{budget} and are genuine exhaustion. \Cref{app:relation-cases} gives all
thirteen with their scores.

\subsection{Verifier Rejection and Acceptance Errors}\label{sec:verifier-errors}

Before the errors, the correct rejections. \Cref{sec:answer-selection} narrates
one in full: the verifier rejected a claim the drafter had fabricated about the
wrong actor. The layer worked exactly as designed. The $52$ questions on which
the verifier fired \textsf{unsupported} are not a population of mistakes.

The two error polarities are, however, very unevenly evidenced --- five cases
against one. That imbalance is not a fact about the verifier. It is a fact about
the instrumentation. It is the finding of this section.

\paragraph{Defining the Error Polarities}\label{sec:polarities}

The verifier's positive class is ``the claim is supported''. A \emph{false
positive} is therefore a claim wrongly \textbf{accepted}, and a \emph{false
negative} a claim wrongly \textbf{rejected}. Because the verifier's own outputs
are labelled \textsf{supported} and \textsf{unsupported}, ``positive'' has no
stable referent for a reader. The rest of this thesis uses the plain-English
terms.

\paragraph{Wrongly Rejected, and Wrongly Accepted}

Five cases are wrongly rejected, all sharing one mechanism and all carrying the
subtype \textsf{structural\_not\_semantic}: the claim is true and the answer
correct. But no single triple states it in the phrasing the drafter used. So a
transitively-true or differently-anchored fact is rejected, and the retry loses
the answer. Three of the five had already resolved the exact gold answer before
the rejection forced the retry. This is the limitation
\Cref{sec:verify-failure-modes} describes by construction, now seen at census
scale. \Cref{sec:future} names the semantic-level check that would address it.

Exactly one clean specimen of wrongful acceptance exists in the whole census.
\textsf{WebQTrn-1597\_8adc06ba} asserts a Seth MacFarlane connection to
\emph{ThunderCats}. The verifier returns \textsf{grounded}. A direct Cypher
query confirms that no such edge exists anywhere in the environment. Its
instructive companion looks like a second defect and is not one.
\textsf{WebQTest-1620} asks when President Wilson was in office. The environment
exposes inauguration events but no date literal. So the system answers with two
inauguration mentions whose entities exist and whose edges are real. The
structural check passed \emph{correctly} while the answer is wrong. That
contrast is the most useful point in this section: structural grounding is a
claim about edge existence, not about answer adequacy. So a verifier that checks
whether a claim is \emph{in} the graph cannot, by construction, check whether it
is \emph{the answer}. \Cref{tab:groundedness}'s Tier-2 column measures that gap.
\Cref{app:verifier-cases} gives all seven cases with their traces.

\paragraph{A Rate for One Polarity, and a Blank for the Other}

Three quantities are easy to conflate. Only one is the denominator this polarity
needs. The verifier is invoked $828$ times over the $800$ questions, since a
question that repairs is verified again afterwards. Its \emph{first} verdict is
\textsf{unsupported} on $52$ questions --- $16$ of $400$ on WebQSP and $36$ of
$400$ on CWQ --- of which $13$ end \textsf{grounded} after a repair, leaving
$761$ questions answered under a grounded verdict and $39$ under an unsupported
one. The five wrongly-rejected cases are drawn from the $52$. So that polarity
has a denominator in the question unit: roughly one rejected question in ten
cost a correct answer. Even that is a lower bound, since a wrongful rejection
that did not change the final answer never entered a census that reads only
failures. The units matter because the run record stores the \emph{first}
verifier entry. So a question repaired into groundedness is filed under
\textsf{unsupported} in the committed \textsf{verifier\_outcome} field: reading
that field as the outcome would put $13$ questions on the wrong side of the
ledger. Treating the $52$ as firings would credit the verifier with $52$ of
$828$ rather than the $52$ of $800$ the rate is computed against.

Wrongful acceptance has \emph{no logged population at all}. Accepted claims are
never persisted (\Cref{sec:instrumentation}). So the wrongly-accepted cases sit
undifferentiated inside the $761$ questions answered under a grounded verdict,
with nothing marking them. The single specimen surfaced only because the
ablation-discordance reading happened to walk that trajectory by hand. What can
be given is the size of the space they hide in, since a blank invites the
reading that the space is small. The verifier accepted $2{,}008$ of the
$2{,}110$ claims it decomposed on test. Neither structural route matches on the
relation or the direction. So every acceptance that did not come from the
entailment check was made by a relation-blind test ---
\Cref{sec:structural-check} shows the record cannot say how many that is beyond
the interval $[39, 2{,}008]$. The two polarities are therefore reported
asymmetrically and deliberately: a rate for rejection, an explicit blank for
acceptance, and no average of the two, since a combined ``verifier error rate''
computed from these numbers would be an artefact of what the logs happen to
record. \Cref{sec:threats} treats this as the instrumentation gap it is.
\Cref{sec:future} gives the one-line fix.

\section{Graph Gaps, the Echo Attractor, and Benchmark
Defects}\label{sec:failures-outside}

The categories in this section belong to no single component. In the first, the
environment cannot express what the question asks. In the second, the graph
offers a true fact one node away from the answer, and the system takes it. In
the third, the benchmark's own label is wrong or its question is
underdetermined.

\subsection{Knowledge-Graph Gaps: Literals, Temporal Qualifiers, and
Ordinals}\label{sec:kg-gaps}

\textsf{kg\_gap} is the third-largest category at $43$ and CWQ's single largest
hedge category at $24$. It invokes the unanswerable-in-environment class defined
in \Cref{sec:unanswerable}. Its five subtypes cluster by mechanism. Those five
account for $39$ of the $43$; \Cref{app:kg-gap-cases} gives them with their
specimens. The remaining $4$ carry no subtype, because the environment gap in
each was specific enough that grouping it with anything else would have been an
invention. They stay in the category total throughout rather than being dropped
silently.

One subtype differs in kind from the other four, and the difference matters for
what follows from them. \textsf{data\_error} is a gap in this particular import.
\textsf{ordinal}, \textsf{numeric\_literal}, \textsf{date\_literal}, and
\textsf{temporal\_qualifier} are gaps in what a triple-and-traversal
architecture can express at all. No larger graph fixes them; they need an
operator. That is why \Cref{sec:future} lists literal and ordinal support with
the architectural items rather than with the repairs.

\subsection{The Echo Attractor: Answering with the Topic or an Intermediate
Entity}\label{sec:echo}

\textsf{echo} names a mechanism that recurs independently of the two mechanical
categories above. The system resolves a real, well-grounded entity that
sits\index{echo attractor} \emph{near} the correct answer in the graph. The
verifier passes it, because it genuinely is grounded. The fact is true. It is
simply not what was asked.

Three shapes appear, and the boundary between them is not always crisp.
\textsf{topic} ($6$ cases) is confusion between entities that share a name or a
tight cluster. \textsf{WebQTest-1707} asks who plays Lex Luthor on
\emph{Smallville} and answers ``John Glover plays Lionel Luthor'', a different
character in the same fictional cluster. Lionel, Alexander, and Lucas Luthor
variants all surface in the trace before the run gives up. \textsf{granularity}
($7$ cases) is the right entity family at the wrong level: \textsf{WebQTest-86}
answers ``Kingdom of Denmark'' where gold is ``Denmark'' --- a genuinely
distinct node, the formal sovereign-state entity rather than the common-name
one. \textsf{WebQTrn-1405\_ced20d84} asks which school founded in 1636 President
Kennedy attended, has \textsf{Harvard College} in its very first candidate list,
and narrows to \textsf{Harvard University} instead --- whose founding-date claim
the verifier then correctly rejects. \textsf{intermediate} --- stopping at a hop
along the way rather than continuing to the target --- has no rows of its own in
this census. The reading pass filed those cases under \textsf{granularity}. The
two are best treated as one family.

The throughline is what makes this a named mechanism rather than a residue
category. None of these are cases where the system knew nothing. Every one
resolves a real, verifiably-true fact through a real edge. The failure is
entirely in \emph{which} true fact gets surfaced --- architecturally different
from picking a wrong edge and from there being no edge at all. The right edge
exists, is used, and lands one node away.

That is also why the echo attractor is invisible to any evaluation treating
systems as independent. \Cref{sec:benchmark-defects} shows all five systems of
the main matrix falling into it together, on the same questions.

\subsection{Benchmark Defects: Gold Noise and Ambiguous
Questions}\label{sec:benchmark-defects}

Two numbers must be kept apart. The consensus pass of
\Cref{sec:gold-quality}\index{benchmark defect}\index{gold noise} \emph{flagged}
$58$ WebQSP and $47$ CWQ questions where at least three of the five systems
agreed on the same non-gold answer. Adjudication confirmed $22$ and $19$ as
genuine defects. The gap between those figures is not measurement noise --- it
is the echo attractor of \Cref{sec:echo}, operating across systems.

Read what the flagged consensus answers have in common. Five systems answer
``rick scott'' where gold lists his professions, ``thor heyerdahl'' where gold
lists his occupations, ``amelia earhart'' where gold is Pilot and Writer --- the
topic entity returned as the answer. Five answer ``belgium'' where gold is the
containing region, ``dominican republic'' where gold is Greater Antilles, and
``san francisco giants'' where gold is the 2014 World Series: each is the
composition's intermediate hop rather than its final one. Five answer
``maryland'' for counties, ``nevada'' for Las Vegas, and ``united arab
emirates'' for Dubai: each is one level too coarse.

\emph{When systems agree with one another against gold, the usual cause is a
shared attractor, not a label error}. That is a finding about evaluation
methodology as much as it is about these systems. A policy of rescoring whenever
three systems agree --- which is a natural thing to want --- would have silently
converted this system's most characteristic failure mode into apparent
correctness. The manual adjudication step is what stands between the two. It is
the reason the pass is reported with both numbers rather than one.

\paragraph{The Confirmed Defects}

\Cref{app:defects} reads the confirmed defects case by case --- the twelve
WebQSP and seventeen CWQ questions whose gold is wrong, the ten and two whose
questions are underdetermined, and where each came from. Two patterns account
for most of the CWQ set, both properties of template generation rather than of
annotation care: a gold list flattened from a multi-valued relation, and an
annotation query that traversed one relation too far.

\section{Discussion}\label{sec:discussion}

\subsection{Findings Against RQ1, RQ2, and RQ3}\label{sec:findings}

\textbf{RQ1 --- does agentic navigation improve multi-hop factual accuracy?}
Yes. The multi-hop qualifier is doing real work. AGR leads every accuracy cell
of \Cref{tab:main}. All eight paired comparisons reject, including both against
the agentic baseline running on the same graph with the same model under the
same budget. The stratum profile is the substance of the answer: on CWQ, AGR's
accuracy \emph{rises} with hop count ($0.46 \to 0.55 \to 0.57$), the only system
that ends above where it started, while Vector-RAG collapses from $0.33$ to
$0.10$ on a context that cannot hold a chain at any radius.

Three boundaries are reported with it. Think-on-Graph hedges marginally less on
CWQ, $22.5\%$ against $23.0\%$ --- the one cell of \Cref{tab:main} AGR does not
lead, and the reason the claim is about the accuracy columns rather than the
whole table. It is also better on CWQ's one-hop stratum, as Vector-RAG is on
WebQSP's. The margin over it is located in the budget: where the shared cap lets
it finish it is marginally ahead on both datasets. So the claim defended here is
that guided search completes under a realistic budget where beam search does not
(\Cref{sec:tog-split}). And that comparison is against a systematically narrower
baseline, since Think-on-Graph prunes to the published algorithm's $40$
relations and $20$ neighbours against AGR's $300$ and $200$, binding on roughly
a third of its expansions and almost none of AGR's (\Cref{sec:baseline-tog}).
The margin is a lower bound on what an equal-width reimplementation would show.
The static graph baseline's per-stratum decay is not offered as evidence here.
Its retrieval radius confounds it (\Cref{sec:baseline-graphrag}).

\textbf{RQ2 --- what does pre-generation verification contribute beyond
navigation?} The answer has three parts that a yes-or-no question would have
merged, which is why it was not asked as one.

\emph{First, navigation eliminates structural hallucination. The verifier is not
what does it.} Both navigating systems assert zero ungrounded entities on WebQSP
--- over $1{,}235$ assertions for AGR and $832$ for Think-on-Graph --- against
$27.1\%$ for parametric answering on the same set. Any system that copies its
answer entities out of triples it traversed inherits this. So attributing it to
the verification layer would be wrong. \Cref{tab:groundedness} is the evidence
that it would be wrong.

\emph{Second, the precision advantage belongs to the architecture. This
experiment cannot award any of it to the layer.} AGR's assertion precision
exceeds Think-on-Graph's by $11.8$ points on CWQ ($67.9\%$ against $56.1\%$) and
its per-question precision by $0.13$ and $0.08$, at a lower hedge rate on
WebQSP. So it is not bought with abstention. That margin is real. What it is not
is evidence for the verification layer. The one condition that isolates the
layer does not reproduce it. Removing the verifier moves per-question precision
by $+0.001$ on WebQSP and $-0.010$ on CWQ (\Cref{tab:ablation}), a sign flip an
order of magnitude below the gap it would have to account for. The mechanism
cannot carry the attribution either: entity filtering runs only on the give-up
path (\Cref{sec:entity-filtering}), which $10$ of $400$ WebQSP and $29$ of $400$
CWQ questions entered. A filter touching $2.5\%$ of a test set cannot produce a
$0.13$ precision margin on it. Two differences the ablation holds fixed remain
live as explanations. Neither is the claim check --- the groundedness filter on
evaluator resolutions (\Cref{sec:groundedness-filter}), which is navigation
rather than verification and runs in every ablation condition, and the drafting
prompt's anti-vacuity and grounding provisions, which the baselines' plain
prompt lacks (\Cref{sec:baseline-constants}). Some share of the precision margin
over the agentic baseline may be prompt rather than structure. This work cannot
say how large that share is.

\emph{Third, verification shows no measurable Hits@1 effect.} Removing it
changes two questions out of $398$. That is not a failure of the layer. It is a
statement about what Hits@1 measures: a metric satisfied by one correct entity
among several cannot see a mechanism that works on the others. What the layer
offers is an answer that arrives with the triples supporting it
(\Cref{sec:output-contract}). That is an auditability claim, not an accuracy
claim. The honest formulation is that \emph{claim verification does not change
how often the system is right; it changes how often the system is right for a
reason it can show}.

\textbf{RQ3 --- which components contribute what, at what token cost?} One
condition of four reaches significance. The planner's effect depends on the
stratum: removing it improves WebQSP by $0.083$ F1 at $p = 0.006$, saving $31\%$
of tokens and $35\%$ of calls, and costs CWQ $0.047$ F1 at $p = 0.088$. That
asymmetry, not a pooled average, is the result. It says how decomposition should
be \emph{applied}, not whether it works. Backtracking, verification, and the
$\alpha$-blend show no detectable effect at $n \approx 200$, each for a
different reason, and \Cref{sec:ablations} separates the reasons out. On cost,
the whole apparatus runs at half the language-model calls of the agentic
baseline and never exhausts the call budget that clips the baseline on $44\%$ of
CWQ questions.

The three answers do not point the same way. The chapter does not force them to.
Navigation is where the accuracy comes from; verification is where the
auditability comes from, the precision margin belonging to the architecture as a
whole rather than to any one component this ablation could isolate; and
decomposition, the component the literature treats as straightforwardly
beneficial, is the one component this work can prove is sometimes harmful.

\subsection{What Agency Buys, and What It Costs}

Agency buys the deep tail: on CWQ the accuracy of the agentic system
\emph{rises} across hop strata while every static system decays. It buys that at
roughly an order of magnitude in tokens over one-shot retrieval. Against the
other agentic system, it buys the same accuracy for half the language-model
calls, never exhausting the call budget that clips the baseline on $44\%$ of
multi-hop questions --- the difference between a guided search and a blind one,
priced. What it costs, beyond tokens, is a longer failure surface. Every
mechanism in this census except \textsf{kg\_gap} and the benchmark defects is a
failure the static baselines cannot have. They have no plan to mis-decompose, no
frontier to abandon, and no draft to check. \textsf{decomposition\_error} at
$38$ cases is the clearest instance: a component added to help contributed a
failure family of its own. \Cref{sec:ablations} shows it is a net negative on
the shallower benchmark. Agency converts a retrieval problem into a control
problem. Control problems fail in more ways.

\subsection{When Verification Helps and When It
Over-Hedges}\label{sec:verification-tradeoff}

The verification layer's contribution is auditability. Its cost is recall. Both
directions show up in this census at case level, and case level is where they
are established, because the aggregate effect of removing the layer is
indistinguishable from zero (\Cref{sec:findings}). The layer helps when a draft
asserts something the traversal never supported. Two cases show it: the De Niro
case of \Cref{sec:answer-selection}, where the verifier rejected a fabricated
actor before it could be emitted, and the Vicksburg-president case of
\Cref{sec:benchmark-defects}, where declining to assert was the honest thing to
do and the metric scored it as a loss. The layer over-hedges when a true claim
has no single edge stating it in the drafted phrasing. The five wrongly-rejected
cases of \Cref{sec:verifier-errors} show that. Three of them had already
resolved the exact gold answer before the rejection forced a retry that lost it.
It is the same trade the development set measured in \Cref{sec:dev-tuning}, and
the same one the ablation could not detect at $n \approx 200$. So the claim is
narrow: verification does not make this system more accurate. It makes the
answers it does give attributable, at a small and measurable cost in recall. It
provides no protection at all against the failure mode this census finds most
often --- an answer that is grounded, connected, true, and not what was asked.

\section{Threats to Validity}\label{sec:threats}

Eleven threats follow. Each is stated against the measurement it bears on, with
the mitigation where one exists.

\textbf{The environment is friendlier than the knowledge base it stands in for.
This bounds every absolute number.} It is the union of per-question subgraphs
the RoG distribution pre-extracted \emph{so as to contain their own question's
answer} (\Cref{sec:source-data}). Taking the union restores distractor mass but
cannot restore what was never extracted. That gap is why answer reachability is
certified at $97.3\%$ and $99.7\%$ (\Cref{sec:validation-gate}) where a BFS
extraction from full Freebase under the same hop bound would not reach those
rates. Accuracies measured here are therefore \emph{not} comparable to
evaluations over full Freebase. The ceilings should be read as part of the
result. The mitigation is real but narrow: every system navigates the identical
environment. So the \emph{relative} ordering is sound even where the absolute
level is optimistic. Every claim in \Cref{sec:findings} is relative for this
reason.

\textbf{An identical environment is not identical access to it. One baseline
sees less of it.} The agentic baseline truncates candidate sets at $40$
relations and $20$ neighbours against this system's $300$ and $200$. Those are
the reimplemented algorithm's own pruning widths. But the narrower cut binds on
about a third of its expansions against effectively none of this system's
(\Cref{sec:baseline-tog}). The comparison holds the graph constant and does not
hold the view of it constant.

\textbf{Every reported number comes from one run per condition. Run-to-run
variance is unmeasured.} The confidence intervals and paired tests resample
questions. So they bound the sampling variance of the question draw and nothing
else. Each system's answer set enters them as fixed. The backbone is not
deterministic at the trajectory level even at temperature $0.0$ ---
\Cref{sec:backbone} measures the frozen model at $8$ of $12$ probes, about one
trajectory in three diverging on a rerun. A divergent trajectory can change an
answer, a hedge, a token count and a call count at once. That figure is itself
only partly auditable. Only the first of four rounds survives as a committed
artifact. A reader who declines to take the other three on trust should
therefore read the stability figure as unquantified rather than as $67\%$. That
substitution makes this threat larger, not smaller. Two things make a rank
reversal from decoding noise alone unlikely: the margins defended here are large
relative to their intervals, and the paired tests are computed per question. But
unlikely is the strongest statement available. The experiment that would bound
the variance was not run, and at \$11.13 against a \$15--20 ceiling it was
affordable. So this is a gap in the design, not a budget constraint. A repeated
matrix reporting across-run spread is the cheapest correction
(\Cref{sec:future}).

\textbf{Entities are keyed by name, so homonyms are merged.} Distinct real-world
entities that share a surface form collapse into one node. Two consequences
follow, in opposite directions, and neither is measured here. A merged node can
manufacture a path between things that are not connected, which inflates
reachability and lets a wrong answer verify as grounded. The same merge can
destroy a distinction the question depends on. Name-keying carries this
inherently; it is not a defect of the construction. But it means structural
groundedness is a claim about \emph{this} graph's node identity, not about the
world.

\textbf{The instrumentation gap is the most serious of the remaining items.} The
verifier logs only the claims it rejects, never the ones it accepts or the
triples that matched them. So wrongful acceptance does not diagnose itself. The
single specimen in this census surfaced only through manual trace
reconstruction, and nothing says how many others sit inside the $761$ grounded
questions. One polarity is reported at census scale and the other anecdotally.
That asymmetry comes from the logging decision, not from the verifier.
\Cref{sec:future} gives the fix.

\textbf{The judge fell short of its pre-specified bar.} Cohen's $\kappa$ came to
$0.6995$ against a threshold of $0.7$ (\Cref{sec:judge}). By the conventional
reading that is substantial agreement, and the miss is $0.0005$. The Tier-2
numbers are still reported as corroborating rather than decisive. The judge also
runs on the same backbone as the systems it judges.

\textbf{Topic entities are given, not linked.} Both graph-navigating systems
seed from the dataset's annotated topic entities. That isolates navigation from
entity linking, as intended. But the reported accuracies assume a linking step
that a deployed system would have to perform itself, with errors of its own.

\textbf{Ablations are underpowered.} Three of four conditions detect nothing at
$n \approx 200$. ``No detectable effect'' is not ``no effect''
(\Cref{sec:ablations}).

\textbf{One relabelling happened after the result it affects was known.} A
development question first classified as mediator-heavy was reclassified as
unanswerable-in-environment once its outcome was in.
\Cref{sec:verification-layer} flags it in place, and the aggregates were not
rescored.

\textbf{The reported accuracy has an unmeasured floor.} The extraction bug of
\Cref{sec:decomposition-errors} turns correct reasoning into a wrong entity list
on one specific question shape, so the measured accuracy on that shape is a
lower bound rather than an estimate.

\textbf{The author adjudicated the benchmark defects.} The exclusion set of
\Cref{sec:gold-quality} rests on single-annotator judgements, made with the
consensus signal visible, and that section records the scope limits. The
accuracy tables are not rescored. The defects run in both directions.

\endinput
